\documentclass[11pt]{article}
\usepackage[T1]{fontenc}
\usepackage[utf8]{inputenc}
\usepackage{mathptmx}
\usepackage{microtype}
\usepackage[margin=2.5cm]{geometry}
\usepackage{amsmath,amssymb,amsthm}
\usepackage{booktabs}
\usepackage{longtable}
\usepackage{array}
\usepackage[hidelinks]{hyperref}
\usepackage{natbib}
\bibliographystyle{plainnat}

\emergencystretch=1em

\newtheorem{proposition}{Proposition}
\newtheorem{theorem}{Theorem}
\newtheorem{corollary}{Corollary}
\newtheorem{lemma}{Lemma}
% Claims are statements about the primitives themselves rather than results
% derived within the model. Kept in a separate series so the distinction is
% visible in the numbering.
\newtheorem{claim}{Claim}

\newcommand{\file}[1]{\href{run:#1}{\texttt{#1}}}
\newcommand{\obj}[2]{\file{#1}\,\texttt{::#2}}
\newcommand{\status}[1]{%
  \ifnum\pdfstrcmp{#1}{PROVED}=0 \textbf{PROVED}%
  \else\ifnum\pdfstrcmp{#1}{VERIFIED}=0 VERIFIED%
  \else\ifnum\pdfstrcmp{#1}{REFUTED}=0 \emph{REFUTED}%
  \else #1\fi\fi\fi}
% Evidence tier, orthogonal to the status word above.
%   L = machine-checked in Lean 4 (kernel-verified, axiom-audited)
%   S = exact symbolic derivation in SymPy (not machine-checked)
%   N = numerical sampling over finitely many economies
\newcommand{\tier}[1]{%
  \ifnum\pdfstrcmp{#1}{L}=0 \textbf{L}%
  \else\ifnum\pdfstrcmp{#1}{S}=0 S%
  \else\ifnum\pdfstrcmp{#1}{N}=0 \emph{N}%
  \else #1\fi\fi\fi}

\title{Threshold Entry in a Fixed-Prize Contest: Proofs}
\author{Anurag Srivastava}
\date{31 August 2026}

\begin{document}
\maketitle

\begin{abstract}
\noindent This note collects the proofs underlying a rebuilt model of
threshold entry in a fixed-prize contest, developed after both referee
reports on a prior submission identified the equilibrium analysis as
unsound (Section~\ref{sec:scope}). Distribution-free throughout: no
distributional family is assumed except where explicitly flagged.
Every claim is tagged \status{PROVED}, \status{VERIFIED} (numerically,
not proved), or \status{REFUTED}, and Table~\ref{tab:index} links each
claim to the file and object that establishes it. Two refuted claims
are the author's own conjectures; one is a referee's. Table~\ref{tab:index}
additionally records an \emph{evidence tier} for each claim --- \tier{L}
machine-checked in Lean~4, \tier{S} exact symbolic derivation, \tier{N}
numerical sampling --- because these are not equivalent grades of
assurance and no claim here is machine-checked end to end from
primitives (Section~\ref{sec:tiers}).
\end{abstract}

\tableofcontents

\section{Scope and verification methodology}
\label{sec:scope}

Verification is split across three tools with an explicit division of
labour. \textbf{SymPy} verifies the analytic content: exact symbolic
integration and algebraic identities, checked against closed-form
function families where the integrals evaluate exactly. \textbf{Lean~4
(core only, no \texttt{sorry})} verifies the discrete and
order-theoretic content: given the analytic sign of a step difference
as an explicit hypothesis, the equilibrium structure (existence,
uniqueness, the prefix property) is machine-checked independently of
the analytic step. \textbf{SciPy/NumPy} supplies numerical evidence
where neither symbolic tool applies: adversarial counterexample
search, sampled comparative statics, and convergence diagnostics.

Mathlib is deliberately not used, so no measure-theoretic content is
formalised in Lean directly; the analytic step enters the Lean
development only as an explicit hypothesis, so the \emph{logical
dependency} of the equilibrium result on the analytic identity is
still machine-checked, even though the identity itself is verified in
SymPy rather than Lean. The reason for staying on core Lean is
auditability, not availability --- \S\ref{sec:tiers} states it, and
\file{verify.sh} measures Mathlib's presence on every run rather than
asserting it either way.

All source files are included in the accompanying bundle. Run
\texttt{./verify.sh} to reproduce every claim in Table~\ref{tab:index}; it
regenerates \file{VERIFICATION.md} from the actual run, including an
axiom audit covering every theorem declared in the Lean file, so that
status file cannot drift from the suites it reports. This document is the
canonical statement of the results; \file{VERIFICATION.md} records only
what the checks printed.

A second and separate suite, \texttt{lit/verify\_lit.sh}, covers the
\emph{positioning} claims --- the assertions this document makes about what
other papers prove and how their results relate to these. It is deliberately
not folded into \texttt{./verify.sh}, because a failure there would mean a
misreading of the literature rather than a defect in a result. Two positioning
statements below rest on machine-checked artifacts in that suite: the
dispersive-order bridge of \S\ref{prop:P9gen} and the kernel correspondence of
\S\ref{prop:PMU}. Those artifacts verify \emph{our reading} of the cited
sources; they do not reprove anything in those sources.

\section{What the three evidence tiers mean}
\label{sec:tiers}

The verification suite has three legs and they do not carry equal weight.
Table~\ref{tab:index} records which apply to each claim.

\paragraph{\tier{L} --- Lean 4.} Kernel-verified. \file{lean/EntryContest.lean}
compiles with no \texttt{sorry} and no user-declared axioms; \file{verify.sh}
runs \texttt{\#print axioms} on every declared theorem and fails if any is
unaudited. Of the 52 theorems, 36 depend on no axioms at all and 16 depend
only on \texttt{propext} and \texttt{Quot.sound}, both Lean core axioms,
pulled in by \texttt{omega} on \texttt{Int} arithmetic. These three counts are
cross-checked against the generated \file{VERIFICATION.md} by
\file{checks/verify\_docs.py}, so they cannot drift from the audit as the
development grows. Verified on Lean 4.33.1, most recently on 2~September 2026;
the file also compiled unchanged on 4.14.0 in an earlier session, which has not
been retested since.

Two limitations are material and are stated here rather than left implicit.
First, Mathlib is not used, so the development is deliberately confined to
the \emph{discrete and order-theoretic} content; every analytic step enters
as an explicit hypothesis. The word \emph{deliberately} is load-bearing and
was previously doing work it had not earned: the bundle's generated
environment note used to record that Mathlib was unreachable, which was false
on at least one machine the suite is run on. \file{verify.sh} now measures
this on every run rather than asserting it, and \file{VERIFICATION.md} carries
the finding. The reason for staying on core Lean is not availability but
auditability --- it is what allows \texttt{\#print axioms} to certify every
one of the theorems, which a Mathlib-based development could not do at the
same granularity. In particular the antitonicity of $\Delta$ --- the
engine of P5 --- enters as the hypothesis \texttt{step\_nonpos} and is
discharged not in Lean but at tier \tier{S}. Second, the Lean theorems are
conditional deductions: they establish that the stated conclusions follow
from the stated hypotheses, not that the hypotheses hold for the model's
primitives. Both facts are by design, and together they mean that
\emph{no claim in this document is machine-checked from primitives to
conclusion}. What is machine-checked is that the logical skeleton is sound
given the analytic inputs.

\paragraph{\tier{S} --- SymPy.} Exact symbolic computation, not machine-checked
proof. \file{checks/verify\_sympy.py} derives the difference identity (S1--S6),
the P1 representation (S7), the sign of $\int\varphi^2\,dK$ (S8), the
$\mu\to0$ benchmark (S9), the P-MU difference identity and weight shape
(S10--S12) and the two algebraic steps of P7 (S18--S20). These are exact
rational-arithmetic results rather than floating-point samples, but they are
produced by a computer algebra system whose correctness is assumed, not
verified. Most are checked over stated families and are therefore
instantiations of a universally quantified claim, not proofs of it; the
exceptions are S1, S18 and S20, which are abstract in $F,G,C$ or symbolic
in $m$.

\paragraph{\tier{N} --- numerical.} Sampling over finitely many economies.
Strong evidence against a claim (a single violation refutes) and weak
evidence for one. Where a claim rests on \tier{N} alone this is flagged in
the table, because such a claim is supported but not established.

\paragraph{The remaining analytic inputs.} Three analytic facts are supplied
rather than proved. Two enter as explicit hypotheses: the sign of the step
difference (\texttt{step\_nonpos}, discharged at \tier{S} by S1--S6), and the
divergence of $\kappa$ at the support floor (discharged at \tier{S} by S16).
Both are stated as hypotheses in the Lean development and as assumptions in
Section~\ref{sec:primitives}, so the dependency is explicit rather than
hidden.

It is worth being precise about what S1--S6 does and does not leave
family-dependent, because ``discharged at tier \tier{S}'' understates it. The
links carrying the argument are abstract in $F$, $G$, $C$ and $K$: the
algebraic step $H_{m+1}-H_m=-K\varphi$ (S1), the chain rule (S2), the sign of
the integrand (S3), and --- the one that encodes an actual modelling
assumption --- the vanishing of the boundary term when
$\varphi(a)=\varphi(b)=0$ (S5). Of the two remaining, S4 is integration by
parts, a theorem of calculus rather than a fact about this model, checked on
five concrete pairs because SymPy cannot decide it in abstract-function form;
and S6 instantiates the assembled identity on four closed-form families. So
the family-dependence sits in the assembly and in a textbook step, not in the
links that do the economic work.

The third is P7's uniform bound. Its two algebraic steps are now verified at
\tier{S}: the factorisation $F^m-H_m=F^m(1-G^{Q-1-m}C)$ holds for abstract
$F,G,C$ with no family assumed (S18), the bracket is non-negative because
$G^{e}\le1$ on $[0,1]$ and $C\le1$ (S19), and the base integral
$\int_0^1u^m\,du=1/(m{+}1)$ is exact \emph{symbolically in $m$} rather than
over finitely many exponents (S20). What remains assumed is monotonicity of
the integral functional --- the passage from $H_m\le F^m$ pointwise to
$\int H_m\,dF\le\int F^m\,dF$, and from $H_m\ge0$ to $\int H_m\,dG\ge0$.
That is the same functional-monotonicity hypothesis already supplied
explicitly in the Lean development for P2 (\S\ref{prop:P2}), so P7 introduces
no new analytic assumption; it reuses one that was already on the books.

\paragraph{Why P7 is not checked end to end.} The conclusion
$\Delta(m)\le V/(m{+}1)$ could be evaluated exactly over the admissible power
families and compared against the cap. That test was tried and rejected as
uninformative: over those families $\max\,(m{+}1)\Delta(m)\approx0.46$, and
$\approx0.79$ over a widened set, so the strictly tighter caps $1/(m{+}2)$,
$1/(m{+}3)$ and $1/(m{+}4)$ pass on every case as well. A check that cannot
separate the true bound from a tighter false one is not evidence for the true
bound, so the \emph{proof steps} are verified rather than the conclusion. The
numerical block in \file{checks/verify\_p679.py} is retained as corroboration
and is subject to the same caveat.

\paragraph{A caveat on P-MU.} The hump of the weight does not by itself show
that first-order stochastic dominance leaves the sign of the P-MU comparison
open. Two exact witnesses with $F\le G$ and opposite signs show it, and both are
proved in Lean (\S\ref{prop:PMU}). An earlier version of this paragraph rested
the point on S11 and S12, which do not establish it.

\paragraph{Why row BM is in the table at all.} Every other row records a
result derived \emph{within} the model. Row BM does not: it is a statement
about utility functions, upstream of the model entirely, and assumptions do
not otherwise earn rows --- there is none for $V>0$, none for independence of
the draws from wealth, none for common knowledge of the wealth profile, though
that last is load-bearing. The closest precedent, $\kappa\to\infty$ as
$w\downarrow c$, is likewise an assumption discharged at tier~\tier{S} (S16)
and has no row of its own.

BM is listed because it is not the assumption. The assumption is
\eqref{eq:burdenmono}, which sits in \S\ref{sec:primitives} with the others.
What is listed is the claim that the assumption is \emph{strictly} weaker than
the concavity it replaces, and that claim does real work: without it the
\S\ref{sec:contribution} sentence asserting a weaker hypothesis is a rewording
rather than a weakening. It also rests on a single knife-edge construction, so
it is exactly the kind of support the tier column exists to make visible. The
implication in the other direction, concavity $\Rightarrow$
\eqref{eq:burdenmono}, is elementary and is left in the prose, matching how
$\kappa\to\infty$ is handled.

\section{Verification index}

\begin{longtable}{@{}p{1.5cm}p{5.0cm}p{1.5cm}p{0.8cm}p{6.0cm}@{}}
\toprule
\textbf{Claim} & \textbf{Statement (short)} & \textbf{Status} & \textbf{Tier} & \textbf{Verified in} \\
\midrule
\endfirsthead
\toprule
\textbf{Claim} & \textbf{Statement (short)} & \textbf{Status} & \textbf{Tier} & \textbf{Verified in} \\
\midrule
\endhead
\bottomrule
\endfoot
BM \S\ref{claim:bmweak} & Burden-monotonicity is \emph{strictly} weaker than
  concavity: a non-concave $u$ satisfies it &
  \status{PROVED} & \tier{S} &
  \file{checks/verify\_sympy.py} (S22a oscillation identity, abstract in
  $k,c,w$; S22b amplitude exactly zero when $p\mid c$; S22c reduced burden
  strictly decreasing; S22d $u''>0$ exhibited). Control: S23, $p=0.7$ breaks
  \eqref{eq:burdenmono} at $w=1.3$. The implication concavity
  $\Rightarrow$ \eqref{eq:burdenmono} is one line and is proved in
  \S\ref{sec:primitives}; S21 records the identity and instantiates it. This
  row is about the primitives, not a result derived within the model ---
  see \S\ref{sec:tiers}. \\[4pt]
ANON \S\ref{sec:primitives} & Anonymity is delivered by wealth-independent
  draws, and count invariance does \emph{not} extend beyond it &
  \status{REFUTED} (own, outside scope) & \tier{N} &
  \file{checks/verify\_anonymity.py}. A1-2: spread of $\Delta$ within a count
  class is exactly $0$ at $\theta=0$, so \S\ref{prop:anon} is verified on
  induced primitives, not assumed; A1-3 control: order $10^{-1}$ once
  $\theta\ne0$. A1-4c/A1-5: for $\theta<0$ count invariance is false, with a
  displayed two-size witness whose margin exceeds the quadrature bound of
  A1-1 by $>10^{3}$. For $\theta>0$: $0$ violations in $20{,}000$ economies,
  which is support, \emph{not} proof (\S\ref{sec:tiers}, tier \tier{N} for a
  universal claim). Bounds the scope of every result in
  \S\ref{sec:proofs}. \\[4pt]
P1 \S\ref{prop:P1} & $\Delta(m)=V\,\mathbb E[\varphi(M_m)]$ &
  \status{PROVED} & \tier{S} &
  \file{checks/verify\_sympy.py} (S7, exact); \file{checks/run\_all.py} (P1, max err $3\times10^{-6}$) \\[4pt]
P2 \S\ref{prop:P2} & FOSD: $F\le G$ pointwise &
  \status{PROVED} & \tier{L} &
  \file{lean/EntryContest.lean} (\texttt{below\_subset\_of\_pathwise}, \texttt{cdf\_le\_of\_pathwise}); \file{checks/verify\_sympy.py} (S13 pathwise gap, S14 exact CDF ordering); \file{checks/run\_all.py} (P2, 0 violations, 4$\mu\times$3 families) \\[4pt]
P3 \S\ref{prop:P3} & $\Delta(m{+}1)-\Delta(m)=-\tfrac{V}{2}\!\int\!\varphi^2\,dK$ &
  \status{PROVED} & \tier{S} &
  \file{checks/verify\_sympy.py} (S1--S6, exact); \file{checks/run\_all.py} (P3, 4th-order convergence to $8\times10^{-12}$) \\[4pt]
P3-cor \S\ref{cor:monotone} & $\Delta$ strictly antitone in $m$ &
  \status{PROVED} & \tier{L} &
  \file{lean/EntryContest.lean} (\texttt{antitone\_of\_step}, \texttt{strict\_of\_step}); \file{checks/run\_all.py} (P4, 1500 adversarial pairs, 0 violations) \\[4pt]
P5 \S\ref{prop:P5} & Existence \& uniqueness of $k^*$ (prefix property) &
  \status{PROVED} & \tier{L} &
  \file{lean/EntryContest.lean} (3 theorems, listed in \S\ref{prop:P5}); \file{checks/run\_all.py} (P5) \\[4pt]
P5-inv \S\ref{prop:count-inv} & Every pure-strategy equilibrium has count $k^*$; identities pinned under a stated condition &
  \status{PROVED} & \tier{L} &
  \file{lean/EntryContest.lean} (\texttt{count\_le\_kstar}, \texttt{count\_ge\_kstar}, \texttt{equilibrium\_count\_unique}, \texttt{members\_below\_kstar}); \file{checks/verify\_equilibria.py} (E1--E6: 20{,}000 instances, exhaustive enumeration, antitonicity control) \\[4pt]
N4 \S\ref{prop:assort} & Assortative set is always an equilibrium and minimises total $\kappa$; ranking is an unweighted utility sum, not resource efficiency &
  \status{PROVED} & \tier{L} &
  \file{lean/EntryContest.lean} (\texttt{cost\_exchange\_le}, \texttt{totalCost\_cons}, \texttt{surplus\_gap\_is\_cost\_gap}); \file{checks/verify\_equilibria.py} (N4, 6000 economies, 1335 with multiplicity; control 3892; N4-WELFARE, 1344 multiplicity economies) \\[4pt]
P6 \S\ref{prop:P6} & $\mu\to0\Rightarrow\Delta(m)=1/(m{+}2)$, indep.\ of $Q$ &
  \status{PROVED} & \tier{S} &
  \file{checks/verify\_p679.py} (P6 block); \file{checks/verify\_sympy.py} (S9); \file{checks/run\_all.py} (P6) \\[4pt]
R1 \S\ref{sec:R1} & Collapse threshold $\bar Q$ with $k^*=0$ &
  \status{REFUTED} (own) & \tier{N} &
  \file{checks/run\_all.py} (R1) \\[4pt]
P7 \S\ref{prop:P7} & $\Delta(m)\le V/(m{+}1)$; $k^*\le V/\kappa$, unif.\ in $Q$ &
  \status{PROVED} & \tier{S}\,/\,\tier{L} &
  Algebraic steps \tier{S}: \file{checks/verify\_sympy.py} (S18 dominance factorisation, abstract $F,G,C$; S19 $G^{e}\le1$ on $[0,1]$; S20 $\int_0^1u^m\,du=1/(m{+}1)$, symbolic in $m$). Deduction bound~$\Rightarrow$~cap \tier{L}: \file{lean/EntryContest.lean} (\texttt{entry\_index\_bounded}). Integral monotonicity assumed, see \S\ref{sec:tiers}. Corroborated \tier{N}: \file{checks/verify\_p679.py}, \file{checks/run\_all.py}. \\[4pt]
R2 \S\ref{sec:R2} & $k^*$ decreasing in $Q$ for all $\mu>0$ &
  \status{REFUTED} (referee) & \tier{N} &
  \file{checks/run\_all.py} (R2) \\[4pt]
P8 \S\ref{prop:P8} & $k^*$ nondecr.\ in $V$, nonincr.\ in $c$ &
  \status{PROVED} & \tier{L} &
  \file{lean/EntryContest.lean} (\texttt{entry\_monotone}, \texttt{nonentry\_monotone}); \file{checks/verify\_p89.py} (P8a, P8b) \\[4pt]
P-UNIQ \S\ref{lem:crossing} & Single crossing $\Rightarrow$ $k^*$ unique; $\Delta$ invariant to wealth &
  \status{PROVED} & \tier{L} &
  \file{lean/EntryContest.lean} (\texttt{single\_crossing}); \file{checks/verify\_p89.py} (P-UNIQ, 300 economies) \\[4pt]
P9 \S\ref{prop:P9} & MPS sign flips exactly at $\bar w=\bar\mu$ (full effect) &
  \status{PROVED} & \tier{L} &
  \file{lean/EntryContest.lean} (\texttt{linSpread\_isPivotSpread}, reducing P9 to P9-gen; \texttt{linSpread\_monotone}); \file{checks/verify\_sympy.py} (S15a--c); \file{checks/verify\_p679.py} (P9 block); \file{checks/verify\_p89.py} (both branches: 230 + 4000 cases) \\[4pt]
P9-gen \S\ref{prop:P9gen} & Sign rule for arbitrary pivot-spreads (pivot $x_0$, not the mean) &
  \status{PROVED} & \tier{L} &
  \file{lean/EntryContest.lean} (\texttt{cost\_falls\_above\_pivot}, \texttt{cost\_rises\_below\_pivot}, \texttt{order\_preserved}; composed into \texttt{entry\_preserved\_above\_pivot}, \texttt{nonentry\_preserved\_below\_pivot}, \texttt{entry\_set\_grows\_above\_pivot}); \file{checks/verify\_p9gen.py} (4 families, 0 violations; 729 control violations) \\[4pt]
R1 \S\ref{prop:tailprop} & One-sided rank condition on two sorted profiles suffices; pivot-spread is a special case &
  \status{PROVED} & \tier{L} &
  \file{lean/EntryContest.lean} (\texttt{entry\_preserved\_of\_up\_to\_margin}, \texttt{nonentry\_preserved\_of\_beyond\_margin}, \texttt{downFromMargin\_imp\_beyondMargin}, \texttt{pivot\_imp\_upToMargin}, \texttt{pivot\_imp\_downFromMargin}); \file{checks/verify\_p9gen.py} (R1: explicit non-pivot witness; 3000 spreads, 977 multi-crossing, 2276 genuine MPS, 0 violations; R1-SCOPE: MPS alone insufficient, 37/4000; R1-WEAK: margin unconstrained, 4000 cases, control 362; R1-FREE: mean-preservation not assumed, 4000 cases) \\[4pt]
BAND \S\ref{prop:band} & Coverage is exactly $k^*\le L$ or $M<k^*$; every
  margin is covered iff the spread is dispersive; and inside the band the
  displacement signs determine nothing &
  \status{PROVED} & \tier{L}\,/\,\tier{N} &
  Coverage criterion \tier{L}: \file{lean/EntryContest.lean}
  (\texttt{band\_empty\_iff}, \texttt{band\_nonempty\_of\_lt}; both
  axiom-free). Faithfulness of the $L/M$ form and sharpness \tier{N}:
  \file{checks/verify\_tailband.py} (T1, $20{,}000$, $0$ mismatches; T2
  equivalence with single crossing; T3 $20{,}889$ covered instances, $0$
  violations, incl.\ $1982/1982$ where both branches force $k^*$ fixed; T4
  two opposite-direction witnesses inside the band --- a refutation, hence
  established; T5 band frequency $47.2\%$ unstructured vs $0.0\%$ dispersive).
  Resolves open item~\ref{item:tailopen}. \\[4pt]
P9-str \S\ref{prop:P9strict} & Strong spread $\Rightarrow$ $k^*$ falls strictly &
  \status{PROVED} & \tier{L} &
  \file{lean/EntryContest.lean} (\texttt{linSpread\_strict\_anti\_in\_lam}, \texttt{linSpread\_unbounded\_below}, \texttt{entry\_ceiling\_of\_marginal\_exit}, \texttt{strict\_drop\_of\_marginal\_exit}); \file{checks/verify\_sympy.py} (S16 $\kappa\to\infty$, S17 strict monotonicity); \file{checks/verify\_p9gen.py} (P9-STRICT, 400 economies) \\[4pt]
P-MU \S\ref{prop:PMU} & Exact condition for sign of $d\Delta(0)/dQ$; $\Delta(0,\cdot)$ quasi-concave in $Q$ when $\varphi$ is single-peaked &
  \status{PROVED} & \tier{L} &
  \file{lean/mathlib/PMU.lean} (\texttt{pmu\_identity}, \texttt{pmu\_sign\_iff}, \texttt{kernel\_rises}, \texttt{kernel\_falls}, \texttt{pmu\_single\_crossing}, \texttt{pmu\_quasiconcave}, \texttt{pmu\_quasiconcave\_uniform}); \file{lean/mathlib/PMUWitness.lean} (\texttt{fosd\_does\_not\_sign}); \file{checks/verify\_sympy.py} (S10--S12); \file{checks/verify\_p89.py} (P-MU, observed sign) \\[4pt]
Alg & $KF-KG=K(F-G)$ (factorisation used in P3) &
  \status{PROVED} & \tier{L} &
  \file{lean/EntryContest.lean} (\texttt{step\_factor}, \texttt{step\_sign}) \\
\end{longtable}
\label{tab:index}

\section{Primitives}
\label{sec:primitives}

One prize of value $V>0$, \emph{exogenous and fixed} --- decoupled from the
wealth distribution. This decouples the prize from inequality, which
is the fix for the prize/affordability conflation raised in referee
comments (Section~\ref{sec:R2-affordability}).

The prize is allocated by rank rather than purchased, and that is a modelled
feature rather than a convenience. \citet{ColeMailathPostlewaite1992} interpret
status as ``a ranking device that determines how well he or she fares in the
nonmarket sector'', and \citet{ColeMailathPostlewaite1995} observe that when
desirable goods ``are allocated as prizes rather than sold, the standard
welfare theorems regarding the Pareto optimality of the outcomes no longer
apply'' --- a remark in their concluding discussion rather than a theorem, but
the right characterisation of the object. Nothing below inherits their
machinery: by P1 the gain is proportional to $V$, so $V$ enters this model only
as a scale factor, and all that is used of it is that it is fixed and
rank-allocated.

$Q$ challengers plus one incumbent. Each player chooses $e_i\in\{0,1\}$:
pay nominal cost $c$ to acquire social capital. After the entry
decision, draws $r_i,s_i$ independent across players \emph{and independent of
the wealth profile}, from \emph{arbitrary} distributions, $s_i\ge0$. Score
\[
  \sigma_i \;=\; \mu r_i + (1-\mu)s_i e_i.
\]
Highest score takes the prize. Static, one shot.

Independence of the draws from wealth is a substantive modelling commitment and
is where Proposition~\ref{prop:anon} bites: it is what makes $F$ and $G$, and
hence $\Delta$, functions of the entrant \emph{count} alone rather than of
which wealths entered. Economies in which wealth buys ability --- so that
$s_i$ is drawn conditionally on $w_i$ --- are outside the scope of everything
below, and deliberately so: the model isolates the pure affordability channel.

\paragraph{What lies outside that scope, and why it matters.} The exclusion is
not a formality, and the boundary is demonstrated rather than asserted. Let
ability be tilted against wealth by a parameter $\theta$: take
$\Pr(s_i\le x)=x^{a_i}$ on $[0,1]$ with $a_i=e^{\theta z_i}$, where $z_i$ runs
from $+1$ at the richest rank to $-1$ at the poorest, so that $\theta=0$
reproduces the assumption above exactly. Then
(\file{checks/verify\_anonymity.py}, $\Delta$ computed by deterministic
quadrature with the discretisation error bounded so that every reported
violation can be compared against it):

\begin{itemize}
\item At $\theta=0$ the spread of $\Delta$ within a fixed count class is
  \emph{exactly} zero (A1-2): Proposition~\ref{prop:anon} is verified on the
  model's induced primitives, not merely assumed. Away from $\theta=0$ that
  spread is of order $10^{-1}$ (A1-3), so the test is not vacuous.
\item For $\theta>0$ --- ability rising with wealth --- count invariance held
  in every economy sampled. This is \emph{support and not proof}: zero
  violations in a finite sample is tier~\tier{N} evidence for a universal
  claim, in the sense of \S\ref{sec:tiers}.
\item For $\theta<0$ --- ability falling with wealth --- count invariance is
  \textbf{false}. With $Q=4$, $\mu=1/2$, $c=1$, CRRA $\gamma=2$ and
  $\theta=-3/2$, the profile
  $w=(4.2197,\,2.9982,\,1.9339,\,1.8760)$, giving
  $\kappa=(0.0736,\,0.1669,\,0.5537,\,0.6085)$, has exactly two pure-strategy
  equilibria: the \emph{poorest} challenger entering alone, and the two
  \emph{richest} entering together --- counts $1$ and $2$. The binding margin
  exceeds the quadrature error by more than three orders of magnitude, so this
  is a genuine counterexample and not a numerical artefact.
\end{itemize}

The mechanism of the failure is worth stating, because it identifies what
anonymity is really buying. Under a negative tilt the ablest challenger is
also the one facing the highest entry cost, so the question of \emph{how many}
enter and the question of \emph{which} of them enter cease to be separable ---
and it is precisely their separation that Proposition~\ref{prop:count-inv}
turns into a count that every equilibrium shares. Since every comparative
static below is a statement about that count, the assumption is load-bearing
for the whole of Section~\ref{sec:proofs}, not only for
Proposition~\ref{prop:anon}. Positive correlation between wealth and ability
appears to be harmless because it reinforces the cost ordering rather than
opposing it; that direction is not claimed.

\paragraph{Payoffs.} Writing $u$ for the utility of wealth,
\begin{equation}
  U_i \;=\; u\big(w_i - c\,e_i\big) \;+\; V\cdot\mathbf{1}\{i \text{ wins}\},
  \label{eq:payoff}
\end{equation}
with $u$ increasing and satisfying the burden condition stated below. Two
features of \eqref{eq:payoff} do real work and are worth stating rather than
leaving to be inferred from the entry condition. The cost is paid \emph{inside} $u$, so it is felt in utility
units and its burden falls with wealth; the prize is \emph{additively
separable} and enters linearly, so it is valued identically by rich and poor
and carries no wealth effect of its own. Inequality therefore operates through
exactly one channel, the cost side.

This is the \emph{privilege} contest of \citet{SchroyenTreich2016} rather than
their rent-seeking contest, and the difference is not cosmetic: were $V$
money added to the resource, it would sit inside $u$, the two wealth effects
would oppose, and under CARA they would exactly cancel. Keeping $V$ outside $u$
is licensed by what the prize \emph{is} --- a rank-allocated non-market good in
the sense of \citet{ColeMailathPostlewaite1992}, not a transfer that can be
spent.

Induced score CDFs: $G$ for a non-investor ($\mu r$), $F$ for an
investor ($\mu r+(1-\mu)s$). All results are stated in $F$ and $G$
only. Write
\[
  \varphi \;\equiv\; G-F \;\ge\;0.
\]
Regularity: $F,G$ continuous, common lower endpoint, $\varphi$
vanishing at both endpoints of the union of supports.

Wealth heterogeneity among challengers, $w_i\sim\Lambda$, with support
bounded below by $c$ and $\kappa(w)\to\infty$ as $w\downarrow c$. Utility
cost of entry $\kappa(w)=u(w)-u(w-c)$, continuous --- the only channel
through which inequality operates.

\paragraph{Burden-monotonicity.} The assumption on $u$ is
\begin{equation}
  w'>w \quad\Longrightarrow\quad u(w')-u(w'-c) \;<\; u(w)-u(w-c),
  \label{eq:burdenmono}
\end{equation}
that is, $\kappa$ is strictly decreasing in $w$. Note the scale at which this
is stated: it constrains the increments of $u$ over spans of the \emph{fixed}
length $c$, not the curvature of $u$ at a point.

Strict concavity of $u$ is sufficient, and this is how earlier versions of
this document stated the primitive. The implication is immediate from
\[
  \kappa'(w) \;=\; u'(w)-u'(w-c),
\]
which is negative whenever $u'$ is strictly decreasing (S21). It is
\emph{strictly} sufficient: there are non-concave $u$ satisfying
\eqref{eq:burdenmono}, and one is exhibited in Claim~\ref{claim:bmweak}
below. Nothing in this document uses concavity anywhere else. Every proof
that touches the cost side runs on $\kappa$, and the only other appearance of
$u$ is $u'(w-c)>0$ in Proposition~\ref{prop:P8}, which is monotonicity rather
than curvature.

The Lean development was always at this generality: \file{lean/EntryContest.lean}
takes burden-monotonicity as the hypothesis
\texttt{hkap}~$:\forall a\,b,\ a\le b\to\kappa(b)\le\kappa(a)$ directly, and
the symbol $u$ does not occur in the file at all. So the restatement above
closes a gap between the prose and the formal layer rather than changing what
was proved --- the same direction as the tail-condition correction, with the
prose having claimed \emph{less} than had been established.

\begin{claim}[Burden-monotonicity is strictly weaker than concavity]
\label{claim:bmweak}
There exists an increasing $u$, non-concave on every neighbourhood of a set of
points of positive measure, satisfying \eqref{eq:burdenmono} on a non-degenerate
interval. Hence the class of admissible $u$ is strictly larger than the concave
class.
\end{claim}

\begin{proof}
Take $u(x)=\sqrt{x}+\varepsilon\sin(kx)$ with $c=1$, $\varepsilon=3/50$ and
$k=2\pi/p$. For any $k$, $c$ and $w$,
\[
  \sin(kw)-\sin\!\big(k(w-c)\big)
  \;=\; 2\sin\!\Big(\tfrac{kc}{2}\Big)\cos\!\Big(k\big(w-\tfrac{c}{2}\big)\Big),
\]
verified symbolically in $k,c,w$ with no family assumed (S22a). The burden
therefore inherits the oscillation only through the amplitude
$2\varepsilon\sin(kc/2)$, which vanishes exactly when $kc$ is a multiple of
$2\pi$, i.e.\ when the period $p$ divides $c$. At $p\in\{1,\tfrac12,\tfrac13\}$
the amplitude is exactly zero (S22b), so
$\kappa(w)=\sqrt{w}-\sqrt{w-1}$ identically, which is strictly decreasing
(S22c). Yet $u''(x)=-\tfrac14x^{-3/2}-\varepsilon k^2\sin(kx)>0$ at every
point with $\sin(kx)=-1$ in the range, and such points are exhibited (S22d).
\end{proof}

\paragraph{How much room this buys is open, and the example is knife-edge.}
The amplitude $2\varepsilon\sin(kc/2)$ vanishes only on a measure-zero set of
periods. At $p=0.7$, which does not divide $c$, \eqref{eq:burdenmono} fails
outright, and the violation is exhibited at $w=1.3$ (S23). Claim~\ref{claim:bmweak}
establishes that the containment is \emph{proper}; it does not establish that
the additional room is economically substantial, and it should not be read as
evidence that non-concavity is generally harmless. Characterising the
tolerance --- the conjecture being that it is governed by the width of the
convex region relative to $c$, since a burden reading only wealth differences
of size $c$ cannot see a convex stretch much narrower than $c$ --- is open and
is recorded as such in \file{TODO.md}.

\paragraph{Information.} The realised wealth profile
$(w_1,\dots,w_Q)$ is \emph{common knowledge} among the challengers, and
entry decisions are simultaneous. This assumption is load-bearing and was
left implicit in earlier versions of this document. It is what licenses
writing the entry condition on the order statistics $w_{(1)}\le\cdots\le
w_{(Q)}$ as in \eqref{eq:entry}, and hence what makes $k^*$ a
\emph{deterministic} count --- one that is the same in every pure-strategy
equilibrium, by Proposition~\ref{prop:count-inv} below.

It does \emph{not} pin the identities of the entrants, and an earlier version
of this document claimed that it did. It does not: with
$\Delta=(12,10,1)$ and $\kappa=(2,5,8)$ --- admissible, since $\Delta$ is
strictly antitone and $\kappa$ increasing in the wealth ranking --- the
entrant sets $\{1,2\}$, $\{1,3\}$ and $\{2,3\}$ are all equilibria, the last
of these with the \emph{richest} challenger absent, because joining as a third
entrant is unprofitable even for her while the two poorer entrants are each
content facing one rival. What complete information delivers is the count and
the existence of the assortative equilibrium, not uniqueness of the entrant
set.

The alternative --- privately known entry costs --- yields a symmetric
Bayesian threshold in cost space and a \emph{binomial} entrant count, as in
\citet{MorenoWooders2011}; there the number of entrants and the identity of the
entering set are both random, being determined by unobserved cost draws. The
mechanical difference is sharper than the informational one: their equilibrium
threshold is a single number $t$, against which every buyer compares her own
cost, because under private information all buyers face the same expected
utility. Here the comparison value depends on \emph{rank} --- the $j$-th
richest faces $\Delta(j-1)$ --- because she knows how many richer rivals enter
ahead of her. A rank-dependent rule collapses to a flat one only if $\Delta$ is
constant, which Corollary~\ref{cor:monotone} forbids. The two models are genuinely different, and
the difference is exactly this assumption, not the shape of $\kappa$. The
complete-information reading is also what the verification artefacts
implement: both \file{lean/EntryContest.lean} and the Python suites operate
on an explicitly ordered wealth profile.

Under complete information with heterogeneous $\kappa(w_i)$ there is no
symmetry to break, so no mixing is required and no tie-breaking rule is
needed. This is the respect in which the present model differs from the
endogenous-entry literature with identical agents, where the entrant count
is pinned but the identities are not --- resolved there by mixed entry
\citep{LevinSmith1994}, by sequential arrival \citep{FuLu2010}, or by
private random delays \citep{MorganOrzenSefton2012}.

Opponent maximum when $m$ other challengers invest:
$H_m = F^m G^{Q-1-m} C$, where $C\in\{F,G\}$ is the incumbent's CDF.
Gain from investing:
\[
  \Delta(m) \;=\; V\Big[\textstyle\int H_m\,dF - \int H_m\,dG\Big].
\]

\section{Terminology}
\label{sec:terminology}

Every English term used in a claim or proof below is bound here to a symbol
built from the primitives of Section~\ref{sec:primitives}. Terms are listed
only where they carry load in a stated result; where a symbol alone would do,
no term is introduced. Nothing in the table is a new object --- each entry is
a name for something already defined --- so the table can be read as a
translation key rather than as a second set of definitions.

\begin{longtable}{@{}>{\raggedright\arraybackslash}p{0.235\textwidth}>{\raggedright\arraybackslash}p{0.185\textwidth}>{\raggedright\arraybackslash}p{0.52\textwidth}@{}}
\toprule
\textbf{Term} & \textbf{Symbol} & \textbf{Definition} \\
\midrule
\endfirsthead
\toprule
\textbf{Term} & \textbf{Symbol} & \textbf{Definition} \\
\midrule
\endhead
\multicolumn{3}{@{}l}{\emph{Players and actions}}\\[2pt]
challenger & $i\in\{1,\dots,Q\}$ & A wealth-heterogeneous player choosing
  whether to compete. \\
incumbent & $C\in\{F,G\}$ & The additional player, entering exogenously;
  appears only through her score CDF $C$ in $H_m$. \\
invest, enter & $e_i=1$ & Pay $c$ and draw the two-component score. \\
entrant & $i$ with $e_i=1$ & A challenger who invests. \\
outsider & $i$ with $e_i=0$ & A challenger who does not. \\
entrant set & $S\subseteq\{1,\dots,Q\}$ & $S=\{i:e_i=1\}$. \\
prize & $V>0$ & Fixed, exogenous, rank-allocated; enters
  \eqref{eq:payoff} additively, outside $u$. \\[4pt]
\multicolumn{3}{@{}l}{\emph{Wealth and cost}}\\[2pt]
wealth & $w_i$ & Challenger $i$'s resource, with $w_i>c$. \\
wealth profile & $(w_1,\dots,w_Q)$ & The realised vector; common knowledge. \\
wealth distribution & $\Lambda$ & The law from which the profile is drawn. \\
rank & $j$ & Index into the descending order statistics
  $w_{(1)}\ge\cdots\ge w_{(Q)}$; rank $1$ is richest. \\
entry cost (utility) & $\kappa(w)=u(w)-u(w-c)$ & The burden of $c$ in
  utility units. \\
burden-monotonicity & $\kappa$ strictly decreasing in $w$ & \eqref{eq:burdenmono}:
  the assumption on $u$, stated at the scale of $c$ rather than pointwise.
  Implied by strict concavity, strictly weaker than it
  (Claim~\ref{claim:bmweak}). The hypothesis \texttt{hkap} in the Lean
  development. \\
entry cost (resource) & $c$ & The nominal cost, identical for all agents.
  Distinct from $\kappa$; the two are never interchangeable. \\
affordability channel & $\kappa$ & The sole route by which the wealth
  distribution affects behaviour, given \eqref{eq:payoff}. \\[4pt]
\multicolumn{3}{@{}l}{\emph{Scores and gains}}\\[2pt]
score & $\sigma_i=\mu r_i+(1-\mu)s_ie_i$ & What the prize is awarded on. \\
score CDFs & $F$, $G$ & Law of $\sigma$ for an entrant ($F$) and an
  outsider ($G$); $F\le G$ pointwise by P2. \\
dominance gap & $\varphi=G-F\ge0$ & The wedge investment buys. \\
opponent maximum & $H_m=F^mG^{Q-1-m}C$ & CDF of the best rival score when
  $m$ other challengers invest. \\
gain from investing & $\Delta(m)=V\!\int\! H_m\,dF-V\!\int\! H_m\,dG$ &
  The right-hand side of the entry condition; antitone in $m$ by
  Corollary~\ref{cor:monotone}. \\
anonymity & $\Delta$ a function of $m$ alone & $\Delta$ depends on the
  \emph{number} of rival entrants, not on which agents they are
  (Proposition~\ref{prop:anon}); requires draws independent of wealth. \\[4pt]
\multicolumn{3}{@{}l}{\emph{Equilibrium objects}}\\[2pt]
entry condition & $\kappa(w_{(j)})\le\Delta(j-1)$ & \eqref{eq:entry}: rank
  $j$ enters when her utility cost is no greater than her gain. \\
single crossing & $\kappa(w_{(j)})-\Delta(j-1)$ increasing in $j$ &
  Lemma~\ref{lem:crossing}; what makes the entry set a prefix. \\
entry count & $k^*$ & The largest $k$ with the entry condition holding at
  every rank $j\le k$. \\
count invariance & $|S|=k^*$ for all equilibria & Every pure-strategy
  equilibrium has the same size
  (Proposition~\ref{prop:count-inv}). \\
identity multiplicity & $S\ne S'$, $|S|=|S'|=k^*$ & Distinct entrant sets
  are equilibria at the same count. \\
assortative set & $\{(1),\dots,(k^*)\}$ & The $k^*$ ranks with the lowest
  $\kappa$, i.e.\ the $k^*$ richest. \\
marginal entrant & rank $k^*$, wealth $w_{(k^*)}$ & The last rank
  satisfying the entry condition. \\
inframarginal entrant & rank $j<k^*$ & An entrant richer than the marginal
  one. \\
outsider rank & rank $j>k^*$ & A rank at which the entry condition fails. \\[4pt]
\multicolumn{3}{@{}l}{\emph{Comparative statics}}\\[2pt]
spread & $w\mapsto w'$ & Any map between wealth profiles; results below
  never assume it preserves the mean unless said. \\
mean-preserving spread & $\sum_j(w'_{(j)}-w_{(j)})=0$ plus
  second-order dominance & The economic case of interest, not a hypothesis
  of Proposition~\ref{prop:tailprop}. \\
displacement & $w'_{(j)}-w_{(j)}$ & The rank-by-rank difference between the
  two sorted profiles. \\
quantile coupling & $\Lambda_\beta^{-1}(r_j)-\Lambda_\alpha^{-1}(r_j)$ &
  The displacement read as a difference of quantile functions; always
  available, since both profiles are sorted. \\
pivot & $x_0$ & The wealth at which the displacement changes sign. \\
pivot-spread & $T$ with $w\ge x_0\Rightarrow Tw\ge w$ and
  $w\le x_0\Rightarrow Tw\le w$ & A spread expanding on both sides of a
  single pivot; the dispersive case. \\
tail condition & $w'_{(j)}\ge w_{(j)}$ for $j\le k^*$; or
  $w'_{(j)}\le w_{(j)}$ for $j>k^*$ & The one-sided rank hypothesis of
  Proposition~\ref{prop:tailprop}; implied by, but weaker than, the pivot
  hypothesis. \\
endogenous margin & $k^*$ moving with the profile & The marginal entrant is
  a different agent before and after; Proposition~\ref{prop:endogenous}. \\[4pt]
\multicolumn{3}{@{}l}{\emph{Aggregates}}\\[2pt]
total entry cost & $\sum_{i\in S}\kappa(w_i)$ & In utility units. Its
  resource counterpart is $|S|c=k^*c$, identical across equilibria. \\
sum of payoffs & $\sum_i U_i$ & Unweighted, hence an interpersonal
  comparison of the cardinal index $u$; see the discussion after
  Proposition~\ref{prop:assort} for what it does and does not entail. \\
selection & a rule picking one equilibrium & Not a refinement: nothing
  claims players coordinate on it. \\
\bottomrule
\end{longtable}

Three pairs are deliberately kept apart because conflating them is what
earlier drafts did. \emph{Entry cost} in utility ($\kappa$) versus in
resources ($c$): only the first varies with wealth, and only the second is
what society gives up. \emph{Count invariance} versus \emph{identity
multiplicity}: the first holds, the second is real, and every comparative
static below is about the count for that reason. \emph{Spread} versus
\emph{tail condition}: the propositions are stated on the second, and the
first is what they are applied to.

\section{Proofs}
\label{sec:proofs}

\subsection{P1: representation}
\label{prop:P1}
\begin{proposition}
$\Delta(m) = V\,\mathbb E[\varphi(M_m)]$.
\end{proposition}
\begin{proof}
By parts, $\int H\,dF = 1-\int F\,dH$ and $\int H\,dG=1-\int G\,dH$, the
boundary terms vanishing since $HF$ and $HG$ run from $0$ to $1$.
Subtracting, $\int H\,dF-\int H\,dG = \int(G-F)\,dH = \mathbb
E[\varphi(M_m)]$.
\end{proof}
Replaces the entire $p_1<1/(Q{+}1)<p_1'$ apparatus of the prior
version with one distribution-free object.

\subsection{P2: first-order stochastic dominance}
\label{prop:P2}
\begin{proposition}
$F\le G$ pointwise, strictly on an interval if $s$ has positive mass
on $(0,\infty)$.
\end{proposition}
\begin{proof}
$X=\mu r+(1-\mu)s\ge \mu r=Y$ pathwise for $s\ge0$, so for every threshold
$t$ the event $\{X\le t\}$ is contained in $\{Y\le t\}$, and monotonicity of
probability under inclusion gives $F(t)\le G(t)$.
\end{proof}

The event inclusion --- the whole probabilistic content of the step --- is
machine-checked as \texttt{below\_subset\_of\_pathwise}, and its consequence
for the CDFs as \texttt{cdf\_le\_of\_pathwise}, with monotonicity of the
probability functional supplied as an explicit hypothesis rather than proved
(no measure theory is available in the Lean development; see
\S\ref{sec:tiers}). The pathwise gap $X-Y=(1-\mu)s$ and the resulting CDF
ordering are checked exactly in \file{checks/verify\_sympy.py} (S13, S14).
Hence $\varphi\ge0$ and, by P1, $\Delta(m)\ge0$: investing never
hurts, with no assumption beyond $s\ge0$.

\subsection{P3: the difference identity}
\label{prop:P3}
\begin{theorem}[key result]
For $K\equiv F^m G^{Q-2-m}C$,
\[
  \Delta(m{+}1)-\Delta(m) \;=\; -\frac{V}{2}\int\varphi^2\,dK \;\le\;0.
\]
\end{theorem}
\begin{proof}
\begin{align*}
\Delta(m{+}1)-\Delta(m)
  &= \int (H_{m+1}-H_m)(dF-dG) = \int K(F-G)(dF-dG) \\
  &= \int K\varphi\,(dG-dF) = \int K\varphi\,d\varphi
   = \int K\,d\!\left(\tfrac{\varphi^2}{2}\right) \\
  &= \Big[\tfrac{K\varphi^2}{2}\Big] - \tfrac12\int\varphi^2\,dK
   = -\tfrac12\int\varphi^2\,dK,
\end{align*}
the boundary term vanishing because $\varphi=0$ at both endpoints. $K$
is a product of CDFs, hence non-decreasing, so $dK\ge0$.
\end{proof}

\subsubsection{Corollary: monotonicity}
\label{cor:monotone}
\begin{corollary}
$\Delta$ is non-increasing in $m$, strictly wherever
$\varphi\not\equiv0$ on $\operatorname{supp}(dK)$.
\end{corollary}

\begin{quotation}
\noindent\textit{This refutes the author's own earlier conjecture.} An
initial sketch argued that because $\varphi$ is hump-shaped (zero at
both ends, positive between), $\mathbb E[\varphi(M_m)]$ need not fall
when $M_m$ rises stochastically, so monotonicity would need a
hazard-rate or single-crossing condition. That reasoning was wrong:
the integration by parts shows the hump shape is irrelevant, because
$\varphi$ enters \emph{squared} against $dK$. Monotonicity is
\emph{unconditional}, and no auxiliary condition is needed.
\end{quotation}

\subsection{P5: equilibrium existence and uniqueness}
\label{prop:P5}
\begin{proposition}
Let $\kappa$ be the marginal entrant's utility cost. Since $\Delta$ is
strictly decreasing (\S\ref{cor:monotone}), the set
$\{m:\Delta(m)\ge\kappa\}$ is a prefix $\{0,\dots,k^*-1\}$, so
\[
  k^* \;=\; \#\{m:\Delta(m)\ge\kappa\}
\]
is the unique pure-strategy equilibrium number of investing
challengers, with $k^*=0$ and $k^*=Q$ as boundary cases.
\end{proposition}
Proved in Lean as \texttt{enters\_downward\_closed},
\texttt{cutoff\_unique}, and \texttt{equilibrium\_is\_threshold}
(\file{lean/EntryContest.lean}), none depending on any axiom.

This delivers the referees' three-regime structure (full withdrawal /
interior / full participation) as a \emph{derived} result rather than
an assumption, and is exactly the asymmetric $k$-of-$Q$ structure
requested in review (\S\ref{sec:R1-asymmetric}).

\subsubsection{P5-inv: the count is an invariant of the game}
\label{prop:count-inv}

Proposition~\ref{prop:P5} constructs the assortative equilibrium and shows the
entry set is a prefix \emph{of that schedule}. Since other equilibria exist
(\S\ref{sec:primitives}), the question is what survives across all of them.
The answer is everything this document goes on to use.

\begin{proposition}[count invariance]
\label{prop:countinv}
Index challengers in increasing order of $\kappa$, equivalently decreasing
wealth. Let $S$ be any pure-strategy equilibrium entrant set: every $i\in S$
satisfies $\kappa_i\le\Delta(|S|-1)$ and every $i\notin S$ satisfies
$\kappa_i>\Delta(|S|)$. Then $|S| = k^*$.
\end{proposition}
\begin{proof}
Write $k=|S|$. Since $S$ has $k$ distinct indices, some member sits at index
$k-1$ or higher, and $\kappa$ is increasing, so
$\kappa_{(k)}\le\Delta(k-1)$; maximality of $k^*$ gives $k\le k^*$.
Conversely, suppose $k<k^*$. Single crossing propagates the entry condition
down from $k^*$, so $\kappa_{(k+1)}\le\Delta(k)$. As $|S|=k$, at least one of
the $k+1$ cheapest agents lies outside $S$, and that agent has
$\kappa\le\Delta(k)$ --- contradicting the outsider condition. Hence
$k = k^*$.
\end{proof}

The two combinatorial steps --- some member at index $\ge k-1$, some index
$\le k$ omitted --- are pigeonhole facts about a $k$-element set of indices;
they enter the Lean development as explicit hypotheses and are discharged by
exhaustive enumeration in \file{checks/verify\_equilibria.py}, which also
confirms the conclusion directly on $20{,}000$ random instances (every
equilibrium of every instance has count $k^*$; identity multiplicity occurs in
$17.3\%$ of them). The order-theoretic core is machine-checked as
\texttt{count\_le\_kstar}, \texttt{count\_ge\_kstar} and
\texttt{equilibrium\_count\_unique}.

\emph{This is what licenses the rest of the document.} P7, P8, P9, P9-gen,
P9-str and P-MU are all statements about $k^*$, and $k^*$ is exactly the
quantity that does not vary across equilibria.

\begin{proposition}[when identities are pinned]
\label{prop:idunique}
If in addition $\kappa_{(k^*+1)}>\Delta(k^*-1)$ --- the first excluded
challenger would decline even the most favourable slot --- then every member of
an equilibrium has index below $k^*$, so the assortative set
$\{1,\dots,k^*\}$ is the unique equilibrium.
\end{proposition}
\begin{proof}
A member satisfies $\kappa_i\le\Delta(k^*-1)$ by
Proposition~\ref{prop:countinv}; if $i>k^*$ then
$\kappa_i\ge\kappa_{(k^*+1)}>\Delta(k^*-1)$, a contradiction. With $k^*$
members all drawn from the first $k^*$ indices, the set is exactly the prefix.
\end{proof}

Machine-checked as \texttt{members\_below\_kstar}, and confirmed without
exception on the $11{,}141$ sampled instances satisfying the condition. Where
it fails, assortativity remains a defensible \emph{selection} --- among
count-$k^*$ equilibria it is the one that fills the slots with the lowest-cost
agents, and so minimises total realised entry cost --- but that is an argument
for a selection rather than a consequence of the primitives, and is treated as
such here.

\subsubsection{N4: assortative entry as a selection}
\label{prop:assortative}

Proposition~\ref{prop:count-inv} fixes the entrant count but not the entrant
set: where $\kappa(w_{(k^*+1)})\le\Delta(k^*-1)$, several distinct sets are
equilibria. The assortative set --- the $k^*$ agents with the lowest entry
cost --- is therefore a \emph{selection} and requires an argument rather than
an assumption. The argument is that it is the cheapest, and that under
anonymity cheapest means efficient.

\begin{proposition}[assortative selection]
\label{prop:assort}
Among the pure-strategy equilibria of the entry game:
\begin{enumerate}
\item the assortative set $\{(1),\dots,(k^*)\}$ is always one of them;
\item it minimises the total entry cost $\sum_{i\in S}\kappa(w_i)$ actually
  paid; and
\item since $\Delta$ depends on the entrant count alone
  (Proposition~\ref{prop:anon}) and that count is the same in every
  equilibrium (Proposition~\ref{prop:count-inv}), the sum of the agents'
  payoffs differs across equilibria \emph{only} through that total. The
  assortative equilibrium therefore maximises the unweighted sum of payoffs.
\end{enumerate}
Part (iii) is a statement about one particular aggregate, and the choice of
that aggregate is a commitment, not a definition; the discussion following
this proposition records what it does and does not entail.
\end{proposition}

\begin{proof}
(i) For $j\le k^*$, single crossing gives $\kappa(w_{(j)})\le
\kappa(w_{(k^*)})\le\Delta(k^*-1)$, so every member is content; for $j>k^*$,
$\kappa(w_{(j)})\ge\kappa(w_{(k^*+1)})>\Delta(k^*)$ by the definition of
$k^*$, so no outsider wishes to join. (ii) All equilibria have $k^*$ members
by Proposition~\ref{prop:count-inv}, and the $k^*$ smallest values of a
finite set minimise the sum over subsets of that size; the assortative set is
exactly the $k^*$ smallest $\kappa$. (iii) Fix two equilibria $S,T$ with
$|S|=|T|=k^*$. The prize term depends on $|S|$ only, so it cancels, leaving
the difference in aggregate payoff equal to
$\sum_{i\in T}\kappa(w_i)-\sum_{i\in S}\kappa(w_i)$.
\end{proof}

The one-step exchange underlying (ii) --- replacing a member by a cheaper
non-member never raises the total --- is machine-checked as
\texttt{cost\_exchange\_le}, with \texttt{totalCost\_cons} for additivity and
\texttt{surplus\_gap\_is\_cost\_gap} for (iii), whose hypothesis
$|S|=|T|$ is where anonymity enters. That the hypothesis is load-bearing and
not decorative is checked directly: across $9022$ pairs of equilibria the
payoff gap equals the cost gap under a gain schedule that is a function of
the count alone, and fails to in $3022$ of the same pairs once the schedule
is allowed to depend on \emph{which} agents entered
(\file{checks/verify\_equilibria.py}, N4 and its control). The full
minimisation in (ii)
follows by iterating exchanges and is the standard sorting fact; it is checked
numerically rather than proved in Lean. Over $6000$ economies, of which
$1335$ exhibit genuine identity multiplicity, the assortative set is an
equilibrium in every case and attains the minimum cost in every case; a
non-assortative count-$k^*$ set is \emph{strictly} costlier in $3892$ control
cases, so the selection is not vacuous
(\file{checks/verify\_equilibria.py}, N4).

\paragraph{What the ranking in (iii) is, and what it is not.}
Three things have to be said, because (iii) is easy to over-read.

First, it carries no resource-efficiency content. Every equilibrium has
$k^*$ members and each pays the same access cost $c$, so the resources
consumed are $k^*c$ in all of them; the equilibria are indistinguishable in
those terms. The ranking exists only because $\kappa$ is a \emph{utility}
increment, $u(w)-u(w-c)$, which differs across agents of different wealth.
Verified directly: over $1344$ economies exhibiting genuine identity
multiplicity, every equilibrium has the same size in every case, while the
total $\kappa$ differs strictly in every case
(\file{checks/verify\_equilibria.py}, N4-WELFARE).

Second, summing $\kappa$ across agents is an interpersonal comparison of the
model's cardinal utility index. That the index is cardinal is already assumed
--- the entry condition \eqref{eq:entry} compares a utility increment against
a prize valued in the same units, which is the privilege-contest form of
\S\ref{sec:primitives} --- so no new cardinality is imported here. What is
new is comparability across agents, and it is a normative commitment rather
than a consequence of the model. Its direction should be stated plainly:
because $\kappa$ falls in wealth, the unweighted sum is smallest when the
wealthiest agents are the ones who enter, and the assortative set is exactly
that set. In all $1344$ multiplicity cases it is the equilibrium whose
members are the richest, strictly (same check). The earlier reading of this
as the choice of ``a planner indifferent to identity'' was wrong: a planner
maximising the unweighted sum is not indifferent to identity, it strictly
prefers wealthy entrants, and would rank the same allocation differently
under any inequality-averse weighting.

Third, it does not compete with the welfare question the referee raised, and
is not an answer to it. This proposition ranks equilibria \emph{at} the
equilibrium count $k^*$; the welfare question is whether $k^*$ is itself too
large, since in a fixed-prize contest entry is rent-dissipating and the
business-stealing logic of \citet{LevinSmith1994} applies. The two operate at
different margins and are logically independent: (iii) selects among the sets
the game admits, and says nothing about whether the game admits too many
entrants. Nothing here is a welfare recommendation.

It remains a selection rather than a strategic refinement in any case:
nothing says players coordinate on the cheapest equilibrium.

\subsection{P6: the \texorpdfstring{$\mu\to0$}{mu to 0} benchmark}
\label{prop:P6}
\begin{proposition}
As $\mu\to0$, $\Delta(m)=\dfrac{1}{m+2}$, independent of $Q$.
\end{proposition}
\begin{proof}
At $\mu=0$, $Y=\mu r\equiv0$, so $G$ is the point mass at $0$ and
$G(x)=1$ for all $x\ge0$. Hence $G^{Q-1-m}\equiv1$ and, with the
incumbent investing, $H_m=F^m\cdot1\cdot F=F^{m+1}$. Then
\[
  \int H_m\,dF = \int_0^1 F^{m+1}\,dF = \frac{1}{m+2},
\]
while $\int H_m\,dG = H_m(0) = F(0)^{m+1}=0$ provided $s>0$ almost
surely. Therefore $\Delta(m)=1/(m+2)$, with no dependence on $Q$.
\end{proof}
Confirms Referee 2's conjecture exactly. The $Q$-independence is
structural: $Q$ enters only through $G^{Q-1-m}$, which collapses to
$1$.

\subsection{R1: no collapse threshold (refuted)}
\label{sec:R1}
The author conjectured a finite $\bar Q$ beyond which $k^*=0$, intended
to replace the old mobility ceiling. \textbf{False.} $\Delta(0)$
converges to a strictly positive limit (numerical evidence in
\file{checks/run\_all.py} (R1)), because non-investor scores are capped
at $\sup\operatorname{supp}(G)=\mu$ while investor scores can exceed
it: $\Delta(m)\to\Pr(X>\mu)\cdot\Pr(\text{beat incumbent})>0$.

\subsection{P7: saturation, uniformly in \texorpdfstring{$Q$}{Q}}
\label{prop:P7}
\begin{theorem}
$\Delta(m)\le V/(m{+}1)$ for every $m$ and every $Q$; hence
\[
  k^* \;\le\; V/\kappa \qquad\text{uniformly in } Q.
\]
\end{theorem}
\begin{proof}
$H_m=F^mG^{Q-1-m}C\le F^m$, since $G\le1$ and $C\le1$. So
\[
  \Delta(m)=\int H_m\,(dF-dG) \;\le\; \int H_m\,dF \;\le\; \int
  F^m\,dF = \frac{V}{m+1}.
\]
Entry at index $m$ requires $\Delta(m)\ge\kappa$, forcing
$m+1\le V/\kappa$. The bound mentions no $Q$.
\end{proof}
Saturation is not asymptotic but a \emph{uniform explicit cap}, which replaces
the prior version's numerical claim. At $\mu\to0$ the bound is nearly tight
against P6's $1/(m{+}2)$.

\paragraph{What is and is not new here.} The accounting behind P7 --- entrant
count times entry cost cannot exceed prize mass --- is not new, and the
proposition should not be read as claiming a new inequality.
\citet{FuJiaoLu2015} obtain the same bound for identical agents with mixed
entry: their rent-accounting inequality~(2) states
$[1-(1-q)^M]V \geq Mq(\Delta + \mathbb{E}[x^\alpha])$, and every feasible point
of their Definition~2 accordingly satisfies $Mq\,\Delta \leq V$, a cap on the
\emph{expected} entrant count. \citet{FuLu2010} reach the same relation as an
equality, their~(7) giving total effort as budget minus $N$ times the entry
cost.

Three things distinguish P7, and the contribution is exactly these. It bounds a
\emph{pure-strategy} count rather than an expected one, so the cap is a genuine
finite index bound rather than a restriction on a probability. It is stated in
\emph{utility} units, $\kappa$ rather than a nominal cost, which is what makes
it interact with the wealth distribution at all. And it holds under
\emph{fixed rules with no designer}: Fu--Lu's $N$ is the count induced by an
optimally chosen contest, whereas here the rules are given and the count is
whatever the wealth ordering produces. Stripped of these, the shared content is
a single monotonicity-and-transitivity step.

The verification splits along the two halves of the argument. The
\emph{implication} --- if $(m{+}1)\Delta(m)\le V$ and challenger $m$ enters,
then $(m{+}1)\kappa\le V$ --- is machine-checked as
\texttt{entry\_index\_bounded} (depends on no axioms), though it is a single
composition of monotonicity with transitivity and carries little content. The
\emph{premise} is where the work is, and its algebraic steps are exact
(\file{checks/verify\_sympy.py}): the dominance factorisation
$F^m-H_m=F^m(1-G^{Q-1-m}C)$ for abstract $F,G,C$ (S18), the bound $G^{e}\le1$
on $[0,1]$ that signs the bracket (S19), and $\int_0^1u^m\,du=1/(m{+}1)$
symbolically in $m$ (S20). The one step not verified is monotonicity of the
integral, which is assumed --- see \S\ref{sec:tiers}, where the reason the
conclusion is \emph{not} tested directly is also recorded. The quadrature
block in \file{checks/verify\_p679.py} (five values of $\mu$, seven of $Q$ up
to $128$, twelve of $m$, no violations, implied cap on $k^*$ observed to
hold) is corroboration rather than the basis of the claim.

\subsection{R2: the \texorpdfstring{$Q$}{Q} comparative static (refuted)}
\label{sec:R2}
Referee 2 conjectured that for $\mu>0$, higher $Q$ shrinks the pie
among investors, so fewer invest. \textbf{Not general}: the sign
depends on $\mu$, reversing at a distribution-dependent crossover
$\mu^*\in[0.1,0.7]$ across the families tested
(\file{checks/run\_all.py} (R2)). Below $\mu^*$, extra non-investors
mainly damage the non-investing deviator, \emph{strengthening} the
incentive to invest as $Q$ grows.

\subsection{P8: prize and affordability, separated}
\label{prop:P8}

Sort challengers by wealth, descending. The $j$-th richest, entering as
the $j$-th entrant, faces $j-1$ others, so her entry condition is
\begin{equation}
  \kappa(w_{(j)}) \;\le\; \Delta(j-1).
  \label{eq:entry}
\end{equation}

\begin{lemma}[single crossing]
\label{lem:crossing}
The left side of \eqref{eq:entry} is non-decreasing in $j$ ($\kappa$ is
decreasing in $w$ and $w_{(j)}$ is decreasing in $j$) and the right side
is non-increasing in $j$ (Corollary~\ref{cor:monotone}). Hence the entry
set is a prefix and $k^*$ is unique.
\end{lemma}

This resolves what was previously listed as an open ``fixed point''
issue: no iteration is required. The reason is structural and worth
stating separately.

\begin{proposition}[anonymity]
\label{prop:anon}
$\Delta(j-1)$ depends on the \emph{number} of entrants only, not on
their wealth: the contest is anonymous, so $F$ and $G$ do not involve
$w$. The wealth distribution therefore enters \eqref{eq:entry} only
through $\kappa(w_{(j)})$.
\end{proposition}

\begin{theorem}[P8]
$k^*$ is non-decreasing in $V$ and non-increasing in $c$.
\end{theorem}
\begin{proof}
By P1, $\Delta(m)=V\,\mathbb E[\varphi(M_m)]$ is proportional to $V$, so
raising $V$ raises the right side of \eqref{eq:entry} pointwise; the
entry set can only grow. For $c$: $\partial\kappa/\partial c =
u'(w-c)>0$, so raising $c$ raises the left side pointwise; the entry set
can only shrink.
\end{proof}

Both directions are instances of a single monotone comparative-statics
lemma, formalised in Lean as \texttt{entry\_monotone} and
\texttt{nonentry\_monotone}; Lemma~\ref{lem:crossing} is
\texttt{single\_crossing} (\file{lean/EntryContest.lean}). None depends
on any axiom.

Note the prize effect applies to rich and poor challengers alike, which
was the referee's objection to the prior version
(\S\ref{sec:R2-affordability}); it is accommodated here rather than
assumed away.

\subsection{P9: inequality and aggregate expenditure}
\label{prop:P9}

\begin{theorem}[P9]
\label{thm:P9}
Let $w\mapsto\bar\mu+\lambda(w-\bar\mu)$, $\lambda>1$,
$\bar\mu=\mathbb E[w]$, be a linear mean-preserving spread. Then:
\begin{enumerate}
\item if the marginal entrant is richer than the mean,
  $w_{(k^*)}>\bar\mu$, then $k^*$ weakly \emph{rises};
\item if the marginal entrant is poorer than the mean,
  $w_{(k^*)}<\bar\mu$, then $k^*$ weakly \emph{falls}.
\end{enumerate}
\end{theorem}
\begin{proof}
Order statistics map as $w_{(j)}\mapsto\bar\mu+\lambda(w_{(j)}-\bar\mu)$,
and
\[
  \big(\bar\mu+\lambda(w_{(j)}-\bar\mu)\big)-w_{(j)}
  \;=\;(\lambda-1)(w_{(j)}-\bar\mu),
\]
which is positive exactly when $w_{(j)}>\bar\mu$. Since $\kappa$ is
decreasing in $w$, the left side of \eqref{eq:entry} falls for entrants
richer than the mean and rises for those poorer. By
Proposition~\ref{prop:anon} the right side $\Delta(j-1)$ is unchanged.
Applying the monotone comparative-statics lemma in each direction gives
the two cases.
\end{proof}

Formally this is P9-gen (\S\ref{prop:P9gen}) instantiated at the pivot
$x_0=\bar\mu$: \texttt{linSpread\_isPivotSpread} machine-checks that
$T w=\bar\mu+\lambda(w-\bar\mu)$ with $\lambda\ge1$ satisfies the
\texttt{PivotSpread} hypothesis, so the two branches follow from
\texttt{entry\_preserved\_above\_pivot} and
\texttt{nonentry\_preserved\_below\_pivot} without a separate argument. The
displacement identity, its sign on each side of the pivot, and mean
preservation are checked symbolically in S15a--c.

\emph{This closes what the previous draft listed as an open item.} The
earlier version proved only the partial effect at fixed $\Delta$ and
flagged the congestion feedback as unresolved. Proposition~\ref{prop:anon}
shows there is no feedback to resolve: $\Delta$ is a function of the
entrant count alone and is invariant to the wealth distribution, so the
partial effect \emph{is} the full effect.

The sign therefore flips exactly at $w_{(k^*)}=\bar\mu$, i.e.\ at entry
rate $1-\Lambda(\bar\mu)$ --- a sharp, distribution-specific prediction
rather than a qualitative ``exclusive versus mass'' statement. Verified on
$\Lambda=\text{lognormal}(0.6,0.45)$: predicted flip at entry rate
$0.4110$, confirmed at all seven threshold quantiles tested. Both branches
are exercised separately in
\file{checks/verify\_p89.py}: 230 cases for the rise branch, 4000 for the
fall branch (629 of them strict), with no violations.

\paragraph{Why this matters.} Inequality \emph{raises} status
expenditure in exclusive competitions and \emph{lowers} it in mass
ones. Non-tautological, since $V$ is fixed by construction, and it
reconciles the motivating empirical pattern with the received result ---
inequality dampens status competition, as in
\citet{ColeMailathPostlewaite1992} (\S IV.A) and, on the intensive margin,
\citet{HopkinsKornienko2004,HopkinsKornienko2010} --- rather than
contradicting it: that finding is recovered as the mass-competition
(below-pivot) case, while exclusive competitions sit in the regime where the
sign flips. The full statement, with precedents and scope, is
\S\ref{sec:contribution}.

\subsection{P9 general: arbitrary pivot-spreads}
\label{prop:P9gen}

The linear spread is a special case. Call $T$ a \emph{pivot-spread about
$x_0$} if $T$ is increasing and
\[
  T(w)\ge w \ \text{ for } w\ge x_0, \qquad
  T(w)\le w \ \text{ for } w\le x_0 .
\]

\begin{theorem}[P9, general form]
\label{thm:P9gen}
Let $T$ be a pivot-spread about $x_0$. If the marginal entrant satisfies
$w_{(k^*)}>x_0$ then $k^*$ weakly rises; if $w_{(k^*)}<x_0$ then $k^*$
weakly falls.
\end{theorem}
\begin{proof}
$T$ increasing preserves the wealth ordering, so the $j$-th order
statistic maps as $w_{(j)}\mapsto T(w_{(j)})$. Since $\kappa$ is
decreasing in $w$, the pivot property gives
$\kappa(T(w_{(j)}))\le\kappa(w_{(j)})$ for $w_{(j)}\ge x_0$ and the
reverse for $w_{(j)}\le x_0$. By Proposition~\ref{prop:anon} the gain
schedule $\Delta(j-1)$ is unchanged. Apply the monotone
comparative-statics lemma in each direction.
\end{proof}

\begin{quotation}
\noindent\textit{Correction to the linear statement.} The relevant
threshold is the \emph{pivot} $x_0$, not the mean. The two coincide for
the linear spread $T(w)=\bar\mu+\lambda(w-\bar\mu)$, which is why the mean
appeared in \S\ref{prop:P9}; a general mean-preserving spread may pivot
elsewhere, and then the flip is at $x_0$. Tested directly: applying the
mean-based rule to non-linear pivot-spreads produces 968 counterexamples,
while the pivot-based rule has none.
\end{quotation}

\paragraph{The pivot-spread class is not ad hoc.} It is, up to a crossing
condition, the standard dispersion order of this literature.
\citet{HopkinsKornienko2009} introduce the \emph{dispersive order} of
\citet{Shaked1982} precisely to sign comparative statics of a competitive
equilibrium under a change in the underlying distribution --- the class of
problem P9-gen belongs to --- and their Definition~1 says $F \leq_d G$ whenever
$G^{-1}(r) - F^{-1}(r)$ is weakly increasing. Transporting this through
$T = G^{-1}\circ F$ makes the displacement $T(w) - w$ weakly increasing, and
then two things follow, both machine-checked in
\file{lit/hopkins\_kornienko/Dispersive.lean}: monotonicity of $T$ is
\emph{implied} rather than assumed (\texttt{dispersive\_imp\_monotone}), and a
zero of the displacement yields the \texttt{PivotSpread} hypothesis about that
point (\texttt{dispersive\_crossing\_imp\_pivot}). So the two hypotheses of
Theorem~\ref{thm:P9gen} are delivered together by a standard order plus a
crossing.

The crossing condition is not decoration. A pure translation $T(w) = w+1$ is
dispersive and monotone yet is a pivot-spread about \emph{no} point at all
(\texttt{shiftUp\_not\_pivotSpread}), so the dispersive order alone does not
suffice.

\paragraph{Precedent.} The closest structural precedent for the pivot rule is
\citet{CostrellLoury2004}'s two-job model, where the wage schedule is governed
by the single sufficient statistic $\hat\mu = F^{-1}(\theta)$, the ability of
the marginal worker between jobs, and they state that the effect of a more
unequal ability distribution ``depends on whether the worker on the margin
between jobs, at quantile $\theta$, lies in the upper or lower tail''. P9-gen
is the entry-count analogue. One difference should be stated rather than
glossed: their marginal quantile $\theta$ is \emph{fixed by technology}, since
jobs come in fixed proportions, so only $\hat\mu$ moves when $F$ does; here
$k^*$ is \emph{endogenous} and is precisely what the comparative static is
about. The analogy is in the sign logic, not in the status of the margin, and
their object is a continuum wage schedule where ours is an integer count.

Their Proposition~10 is also a caution worth carrying: under Cobb--Douglas
crowding, where the assignment re-optimises, the direction of the inequality
effect \emph{reverses} relative to their fixed-proportions Proposition~6. The
analogue of fixed assignment here is that $\Delta$ does not move under a wealth
spread --- and that is not an assumption but Proposition~\ref{prop:anon}. The
literature's own experience is that this sign flips when the allocation
re-optimises; P9 is safe because the absence of feedback was proved, not
assumed.

The hypothesis is not vacuous. Maps that contract near the pivot --- for
instance $T(w)=x_0+\operatorname{sign}(d)|d|^{p}$ with $d=w-x_0$ and
$p>1$, for which $|d|^p<|d|$ when $|d|<1$ --- are increasing but are
\emph{not} pivot-spreads, and they do violate the sign rule (729
violations in the control run). Four qualifying families (linear,
quadratic, mixed, exponential) give zero violations across 1500 economies
each. Verified in \file{checks/verify\_p9gen.py}; the two cost lemmas are
\texttt{cost\_falls\_above\_pivot} and \texttt{cost\_rises\_below\_pivot},
with order preservation as \texttt{order\_preserved}
(\file{lean/EntryContest.lean}, no axioms).

\subsubsection{N1: the margin is endogenous, and the rule survives it}
\label{prop:endog}

Theorem~\ref{thm:P9gen} is stated at the marginal index. That index is not
fixed: $k^*$ is determined by \eqref{eq:entry} and moves when the
distribution moves, so the marginal entrant is in general a different agent
before and after the spread. The rule nevertheless needs only a single
comparison at the \emph{pre}-spread margin.

\begin{proposition}[endogenous margin]
\label{prop:endogenous}
Let $T$ be a monotone pivot-spread about $x_0$, let challengers be sorted
descending, and let $k^*$ be the pre-spread entrant count.
\begin{enumerate}
\item If $w_{(k^*)}>x_0$ then no pre-spread entrant exits, so the post-spread
  count satisfies $k^{*\prime}\ge k^*$.
\item If $w_{(k^*)}<x_0$ then no pre-spread outsider --- that is, no agent at
  a rank $j>k^*$ --- enters, so $k^{*\prime}\le k^*$.
\end{enumerate}
No post-spread margin need be located.
\end{proposition}

\begin{proof}
Descending order does the work. For (i), any $j\le k^*$ has $w_{(j)}\ge
w_{(k^*)}>x_0$, so every inframarginal entrant is also above the pivot; by
the pivot property $T(w_{(j)})\ge w_{(j)}$, so $\kappa$ falls weakly and
\eqref{eq:entry} continues to hold at every such $j$. For (ii), any $j>k^*$
has $w_{(j)}\le w_{(k^*)}<x_0$, so every \emph{outsider} is also below the
pivot, $T(w_{(j)})\le w_{(j)}$, $\kappa$ rises weakly, and \eqref{eq:entry},
which failed at those $j$ by the definition of $k^*$, continues to fail. The
argument is confined to $j>k^*$: at the margin itself \eqref{eq:entry} holds,
so there is nothing there to preserve. Monotonicity of $T$ preserves the
descending order, so Lemma~\ref{lem:crossing} applies to the post-spread
profile and $k^{*\prime}$ is well defined; $\Delta$ is unchanged by
Proposition~\ref{prop:anon}.
\end{proof}

This is stronger than the fixed-index reading of the same rule, and it is the
respect in which the present setting differs from its closest structural
precedent. In \citet{CostrellLoury2004}'s two-job model the marginal quantile
$\theta$ is fixed by technology --- jobs come in fixed proportions, so
$\theta$ does not move when the ability distribution does, and only
$\hat\mu=F^{-1}(\theta)$ shifts. Here the marginal index is itself the object
the comparative static is about, and the sign rule holds anyway.

That difference should not be read as this proposition dominating theirs, and
it is not claimed here. Costrell--Loury sign a smooth aggregate --- a wage
schedule --- under a general spread, which is a harder object than an integer
threshold count and is why their argument needs machinery this one does not.
What is being compared is not two theorems of differing strength but two
formulations: moving the choice to the extensive margin replaces the
aggregate with a threshold statistic, and it is that replacement which makes
an endogenous margin tractable at all. The modelling choice is the
contribution; the proposition is what the choice buys, and it is
correspondingly easy once the choice is made.

The hypothesis is not slack: comparing the pivot against the \emph{richest}
challenger rather than the marginal one produces violations in $154$ of
$2744$ control cases (\file{checks/verify\_p9gen.py}, N1 CONTROL), so the
marginal index is the load-bearing comparison point. Machine-checked as
\texttt{entry\_preserved\_below\_marginal\_index} and
\texttt{nonentry\_preserved\_beyond\_marginal\_index}, with
\texttt{descending\_preserved} for well-definedness; checked numerically with
$k^*$ recomputed on both sides over $20\,000$ economies, $0$ violations, the
count actually moving in $393$ of them.

\subsubsection{R1: a one-sided tail condition, and the spreads it covers}
\label{prop:tail}

Theorem~\ref{thm:P9gen} assumes $T$ is a pivot-spread: the displacement
changes sign exactly once. That is the dispersive case. A general
mean-preserving spread need not be dispersive --- its quantile difference may
cross zero arbitrarily often. The pivot hypothesis is nonetheless stronger
than the argument requires.

Compare the two \emph{sorted} profiles rank by rank. This is always
available: the rank-wise comparison of two distributions is their quantile
coupling, $w'_{(j)}-w_{(j)}=\Lambda_\beta^{-1}(r_j)-\Lambda_\alpha^{-1}(r_j)$,
and it is monotone by construction, so no assumption is needed to obtain it.

\begin{proposition}[tail condition]
\label{prop:tailprop}
Let $w$ and $w'$ be any two wealth profiles, each sorted descending, and let
$k^*$ and $k^{*\prime}$ be the entry counts they induce through
\eqref{eq:entry}.
\begin{enumerate}
\item If $w'_{(j)}\ge w_{(j)}$ for every $j\le k^*$, then $k^{*\prime}\ge k^*$.
\item If $w'_{(j)}\le w_{(j)}$ for every $j>k^*$, then $k^{*\prime}\le k^*$.
\end{enumerate}
Nothing is assumed about the displacement on the other side of the margin,
which may therefore change sign any number of times. In (ii) nothing is
assumed at the margin either: rank $k^*$ is the marginal \emph{entrant}, so
her wealth may move in either direction.
\end{proposition}

The hypotheses are worth reading for what they omit. Two properties are used
and no others: $\kappa$ is non-increasing in wealth, which is concavity and
nothing more (\S\ref{sec:primitives}), and both profiles are sorted, so that
Lemma~\ref{lem:crossing} defines each count. There is no requirement that
$w'$ be obtained from $w$ by a spread, by a mean-preserving transformation,
or by any transformation at all --- in particular
$\sum_j\bigl(w'_{(j)}-w_{(j)}\bigr)=0$ is nowhere used, and the two profiles
need bear no relation to each other beyond the rank-wise inequality assumed.
The Lean development has always been at this generality
(\texttt{entry\_preserved\_of\_up\_to\_margin} quantifies over arbitrary
\texttt{w w'}); what is corrected here is the surrounding statement, which
described the result as being about spreads.

Mean-preserving spreads are the case of economic interest, and everything
below applies the proposition to them. But they are the application, not the
hypothesis, and the distinction matters for what may be claimed: the
proposition is not a comparative-static result about $\Lambda$, it is a
monotonicity result about a threshold statistic, which is then \emph{used}
to obtain comparative statics whenever a spread happens to satisfy its
condition.

\begin{proof}
(i) For $j\le k^*$, $\kappa$ decreasing gives $\kappa(w'_{(j)})\le
\kappa(w_{(j)})\le\Delta(j-1)$, so \eqref{eq:entry} still holds at every such
$j$, and the count cannot fall. (ii) For $j>k^*$ the agent is an outsider, so
$\kappa(w_{(j)})>\Delta(j-1)$ by the definition of $k^*$; the hypothesis
gives $\kappa(w'_{(j)})\ge\kappa(w_{(j)})$, so \eqref{eq:entry} still fails
and the count cannot rise. Note that the argument is available only at
$j>k^*$: at $j=k^*$ the entry condition \emph{holds}, which is what makes
that agent the margin, so no conclusion is drawn there and none is needed.
\end{proof}

The pivot hypothesis implies the tail condition
(\texttt{pivot\_imp\_upToMargin}, \texttt{pivot\_imp\_downFromMargin}), so
Theorem~\ref{thm:P9gen} and Proposition~\ref{prop:endogenous} are special
cases. The converse fails, which is the content: an explicit witness with
$Q=8$, displacement $(0.02,0.02,0.02,0.30,-0.20,0.15,-0.10,\cdot)$ closed to
sum zero, is mean-preserving, keeps the profile sorted, changes sign three
times --- so it is \emph{not} a pivot-spread --- satisfies the tail
condition, and moves $k^*$ from $3$ to $4$. Randomised: $3000$ spreads, of
which $977$ change sign at least twice and $2276$ are mean-preserving spreads
in the second-order-dominance sense rather than merely mean-preserving
perturbations; no violations (\file{checks/verify\_p9gen.py}, R1).
Machine-checked as \texttt{entry\_preserved\_of\_up\_to\_margin} and
\texttt{nonentry\_preserved\_of\_beyond\_margin}, the latter on the
strict-outsider hypothesis stated above;
\texttt{downFromMargin\_imp\_beyondMargin} records that the earlier
formulation, which also signed the displacement at the margin, assumed more
than the conclusion needs. Checked numerically with the margin and every
inframarginal rank left free to move either way: $4000$ cases, $0$
violations, against a control in which no rank is signed and $k^*$ rises in
$362$ of $4000$ (\file{checks/verify\_p9gen.py}, R1-WEAK).

The condition is a genuine hypothesis and not a restatement of
mean-preservation, in both directions. It is not implied by it: among $4000$
sampled spreads that are mean-preserving in the second-order-dominance sense,
$37$ have a quantile difference that turns negative at some rank at or above
the margin, and in those cases $k^*$ \emph{falls} --- witness $Q=10$,
$k^*=8\to7$ (\file{checks/verify\_p9gen.py}, R1-SCOPE). Nor does it imply
mean-preservation: over $4000$ profile pairs satisfying the condition and
deliberately \emph{not} mean-preserving, the count never falls
(\file{checks/verify\_p9gen.py}, R1-FREE), confirming on the checked layer
what the proof already shows --- that
$\sum_j\bigl(w'_{(j)}-w_{(j)}\bigr)=0$ is not among the hypotheses. Which
mean-preserving spreads satisfy the condition is settled in
\S\ref{prop:band} below, and the witness above is what makes that question
non-vacuous.

Why this works where the assignment literature needs more machinery: $k^*$ is
a threshold statistic, not an integral. \citet{Suen2007} and
\citet{CostrellLoury2004} sign smooth aggregates and must therefore control
the quantile difference at every rank, which is what forces integration by
parts and a log-concavity or monotone-weight condition. The entry count reads
the quantile function only down to the margin, so one-sided information
suffices and no such condition is needed.

\subsubsection{The band: which margins the tail condition covers, and why it
cannot be weakened}
\label{prop:band}

Write the rank-by-rank displacement of the two sorted profiles as
$d_j=w'_{(j)}-w_{(j)}$, rank $1$ richest, and set
\[
  L=\min\{j: d_j<0\},\qquad M=\max\{j: d_j>0\},
\]
the richest rank whose wealth falls and the poorest rank whose wealth rises.
The two hypotheses of Proposition~\ref{prop:tailprop} are then exactly
$k^*\le L$ (rise branch) and $M<k^*$ (fall branch), an equivalence verified
over $20\,000$ instances with no mismatch (\file{checks/verify\_tailband.py},
T1). The proposition is therefore silent precisely on margins with
$L<k^*\le M$. Call that range the \emph{band}.

Two things follow, and together they close the question.

\paragraph{What dispersiveness buys, exactly.} Every margin is covered if and
only if $M\le L$ --- that is, if and only if no rank that rises is poorer than
a rank that falls, which is precisely single crossing of the displacement.
This is machine-checked: \obj{lean/EntryContest.lean}{band\_empty\_iff},
axiom-free, together with the contrapositive
\obj{lean/EntryContest.lean}{band\_nonempty\_of\_lt}. So a pivot-spread is not
a convenient special case of the tail condition; it is \emph{the} class on
which the rule holds at every margin, and the band is the exact price of
leaving it.

\paragraph{The hypothesis is sharp.} Inside the band both directions occur.
Over $40\,000$ genuine mean-preserving spreads the rule is never violated
where it applies --- $20\,889$ covered instances, $0$ violations, including
$1\,982$ in which \emph{both} branches apply and the rule accordingly forces
$k^*$ to be unchanged, which it is in every one. Inside the band, by contrast,
$k^*$ rises in $4\,785$ cases and falls in $130$
(\file{checks/verify\_tailband.py}, T3--T4). Two witnesses with the same
displacement-sign class and opposite conclusions are displayed there. Since a
single such pair refutes any purported strengthening, this is established
rather than sampled: \textbf{no conclusion about the direction of $k^*$ can be
drawn from the displacement signs alone once the margin lies in the band.}
The tail condition is not a convenient sufficient condition awaiting
improvement; it is the whole of what the sign pattern supports.

\paragraph{A caution about how large the band is.} It is tempting to quote a
frequency, and misleading to do so. Under the unstructured generator that
produced the R1-SCOPE witness --- displacements drawn independently and then
balanced --- $47.2\%$ of admissible spreads leave the margin in the band. That
number describes the generator, not the economics: multi-crossing is typical
by construction there. Under a linear mean-preserving spread and under a
mean-preserving increase in the dispersion of a lognormal, the figure is
$0.0\%$ in $8\,000$ instances each, since both families are dispersive and
Lemma \obj{lean/EntryContest.lean}{band\_empty\_iff} then applies at every
margin (T5). The honest statement is the conditional one: the band is empty
for dispersive spreads and can be substantial for spreads whose quantile
difference oscillates, and no single frequency should be attached to it.

\paragraph{What remains open.} Only this: whether some statistic of the spread
\emph{other} than the displacement signs --- its magnitudes, or a weighted
functional of the quantile difference of the kind
\citet{CostrellLoury2004} use --- predicts the direction of $k^*$ inside the
band. The present result says the signs do not, and says nothing about
whether something else might.

\subsection{P9 strictness}
\label{prop:P9strict}

\begin{proposition}
\label{prop:strict}
Suppose the marginal entrant is strictly below the pivot,
$w_{(k^*)}<x_0$, and consider the linear family
$T_\lambda(w)=x_0+\lambda(w-x_0)$. Then $k^*(\lambda)$ is non-increasing
in $\lambda$, and there exists $\bar\lambda$ such that
$k^*(\lambda)<k^*(1)$ for all $\lambda>\bar\lambda$.
\end{proposition}
\begin{proof}
Non-increasingness is Theorem~\ref{thm:P9gen} applied at each
$\lambda$. For strictness: $\kappa$ is continuous and strictly
decreasing in $w$, with $\kappa(w)\to\infty$ as $w\downarrow c$. Since
$w_{(k^*)}<x_0$, $T_\lambda(w_{(k^*)})$ decreases without bound in
$\lambda$, so $\kappa(T_\lambda(w_{(k^*)}))$ eventually exceeds the fixed
value $\Delta(k^*-1)$ and the marginal entrant exits. Because the entry
set is a prefix (Lemma~\ref{lem:crossing}), losing the marginal entrant
strictly reduces the count.
\end{proof}

The three steps are separated in the verification. That
$T_\lambda(w_{(k^*)})$ is strictly decreasing in $\lambda$ and unbounded
below is machine-checked (\texttt{linSpread\_strict\_anti\_in\_lam},
\texttt{linSpread\_unbounded\_below}). That losing the marginal entrant
strictly reduces the count --- because the post-spread entry set is a prefix,
so no index at or above the exiting one can enter --- is machine-checked
(\texttt{entry\_ceiling\_of\_marginal\_exit},
\texttt{strict\_drop\_of\_marginal\_exit}). The remaining step is analytic and
is \emph{not} proved in Lean: that $\kappa$ diverges at the support floor, so
that a low enough wealth makes the entry cost exceed the fixed gain
$\Delta(k^*-1)$. It enters as the boundary condition stated in
Section~\ref{sec:primitives} and is checked symbolically for the log and
CRRA families in \file{checks/verify\_sympy.py} (S16, S17). Verified in
addition over 400 qualifying economies: $k^*$ is monotone in $\lambda$ and
strictly lower at $\lambda=100$ than at $\lambda=1$ in every case.

\subsection{P-MU: the \texorpdfstring{$Q$}{Q} comparative static, characterised}
\label{prop:PMU}

\begin{proposition}
For $Q\ge1$, $\Delta(0,Q{+}1)-\Delta(0,Q) = -V\int W\,(dF-dG)$ with
$W \equiv C\,G^{Q-1}(1-G)$, where $C$ is the incumbent's CDF. With $V>0$,
$\Delta(0)$ rises from $Q$ to $Q+1$ if and only if
\[
  \int W\,dF \;<\; \int W\,dG .
\]
\end{proposition}
The identity holds for every incumbent CDF $C$
(\file{lean/mathlib/PMU.lean}, \texttt{pmu\_identity}, \texttt{pmu\_sign\_iff}).
The factor $G^{Q-1}(1-G)$ rises in $G$ up to $(Q-1)/Q$ and falls beyond it
(\texttt{kernel\_rises}, \texttt{kernel\_falls}). First-order stochastic
dominance alone does not sign the comparison. Two pairs with $F\le G$, no atoms
and an investing incumbent give opposite signs at $Q=1$ and at $Q=2$. The pair
$F=x$, $G=1-(1-x)^2$ on $[0,1]$ gives $V/60$ and $V/420$, while the pair
$F=x^2$, $G=x$ gives $-VQ/\big((Q+2)(Q+3)(Q+4)\big)$ at every $Q$
(\file{lean/mathlib/PMUWitness.lean}, \texttt{fosd\_does\_not\_sign},
\texttt{pairB\_step}).

\paragraph{Provenance of the weight.} The factor $G^{Q-1}(1-G)$ has the form of
the log-supermodular kernel $F^{k-1}(1-F)$ that \citet{RyvkinDrugov2020} use for
their comparative statics in the number of players, under $G \leftrightarrow F$
and $Q \leftrightarrow k$. The peak at $(Q-1)/Q$ is our computation, which
\citet{RyvkinDrugov2020} do not state. The extra factor $C$ here is the
incumbent's CDF and carries no $Q$. Their corresponding object is
$b_k = \mathbb{E}[f(X_{(k-1:k-1)})]$, the individual marginal benefit of effort,
which is the right comparator for $\Delta(0)$, whereas their hazard-rate result
concerns aggregate effort and a different order statistic.

\paragraph{The single-crossing step carries over.} Since $f-g=-\varphi'$, the
identity reads $\Delta(0,Q{+}1)-\Delta(0,Q)=V\int G^{Q-1}(1-G)\,C\varphi'\,dx$,
and $C\varphi'$ crosses $+-$ whenever $\varphi$ is single-peaked, the
orientation that the variation-diminishing step of
\citet{RyvkinDrugov2020} needs. An earlier version of this paragraph dropped
the leading minus sign of the identity and concluded that the orientation was
reversed. The step is proved on measures, without densities. If $dF-dG$ is
nonpositive on $(-\infty,x_0]$ and nonnegative on $(x_0,\infty)$, then for
$1\le Q_1\le Q_2$
\[
  \Delta(0,Q_2{+}1)-\Delta(0,Q_2) \;\le\; G(x_0)^{Q_2-Q_1}
  \big[\Delta(0,Q_1{+}1)-\Delta(0,Q_1)\big]
\]
(\texttt{pmu\_single\_crossing}). Once the difference is nonpositive it stays
nonpositive, so $\Delta(0,\cdot)$ is quasi-concave on $Q\ge1$, rising and then
falling, where either phase may be absent (\texttt{pmu\_quasiconcave}). The
hypothesis holds at $x_0=\mu$ when $r$ is uniform on $[0,1]$, for any law of
$s\ge0$ and any $0<\mu\le1$ (\texttt{uniform\_left}, \texttt{uniform\_right},
\texttt{pmu\_quasiconcave\_uniform}). Whether the hypothesis holds for other
laws of $r$ is open.

\section{Contribution and its precedents}
\label{sec:contribution}

\paragraph{The claim.} For a fixed prize awarded by rank, the effect of
resource inequality on \emph{how many people compete} has a determinate sign,
and the sign is set by a single observable: whether the marginal entrant
stands above or below the pivot of the spread. A \emph{dispersive} spread of
the resource distribution --- one whose displacement changes sign exactly
once, the leading case being the linear mean-preserving spread
$w\mapsto\bar\mu+\lambda(w-\bar\mu)$ of Theorem~\ref{thm:P9} ---
\emph{widens} entry when the marginal competitor is richer than the pivot and
\emph{shrinks} it when poorer (Theorem~\ref{thm:P9gen}), strictly so under the
boundary condition of Proposition~\ref{prop:strict}. What the proof actually
uses is weaker still: a one-sided condition on the quantile difference down to
the margin, with mean-preservation nowhere among the hypotheses
(Proposition~\ref{prop:tailprop}). Within that class the rule holds globally,
distribution-free, and under burden-monotonicity of $u$ alone
(\eqref{eq:burdenmono}; strict concavity is sufficient but not necessary,
Claim~\ref{claim:bmweak}).

That class is a genuine restriction, and it belongs with the claim rather than
only with the scope discussion at the end of this section. A
Rothschild--Stiglitz mean-preserving spread need \emph{not} be dispersive: its
quantile difference may cross zero any number of times, and $37$ of $4000$
sampled genuine mean-preserving spreads fail the tail condition with $k^*$
moving against the rule (\S\ref{prop:tail},
\file{checks/verify\_p9gen.py}~R1-SCOPE). So the sign rule is a theorem about
dispersive spreads and about the tail condition, and not about mean-preserving
spreads as such; it should not be quoted in the latter form. Which spreads
satisfy the condition is settled in \S\ref{prop:band}: coverage is exactly
$k^*\le L$ or $M<k^*$ in the notation there, every margin is covered if and
only if the spread is dispersive, and --- the substantive half --- inside the
band the displacement signs determine nothing, so the hypothesis is
\emph{sharp} rather than merely sufficient.

The economic content is a reconciliation. The received result --- inequality
dampens status competition --- goes back to
\citet{ColeMailathPostlewaite1992}, whose \S IV.A example has the more
compact income distribution producing the higher savings rate, and is
developed on the intensive margin by \citet{HopkinsKornienko2004} and
\citet{HopkinsKornienko2010}, where less dispersed endowments intensify
conspicuous consumption. That finding is here recovered as one branch: mass
competitions, in which the marginal entrant sits low in the distribution.
Exclusive competitions, in which the marginal entrant is already wealthy,
sit on the other branch, where inequality \emph{feeds} entry. The received
result is the below-pivot case of a two-branch rule, not the general case.
This is non-tautological because $V$ is fixed by construction
(\S\ref{prop:P8}): nothing about the prize moves when the distribution
spreads.

The margin is endogenous and the rule survives that
(Proposition~\ref{prop:endogenous}): the hypothesis is a single comparison at
the pre-spread margin, and the wealth ordering carries the rest, so no
post-spread margin need be located. \citet{CostrellLoury2004}, the closest
structural precedent, work with a marginal quantile fixed by technology. The
claim here is not that the present result is the stronger one --- they sign a
smooth aggregate under a general spread, which is the harder object --- but
that the extensive-margin formulation replaces that aggregate with a
threshold statistic, and an endogenous margin is tractable in the second
setting and not the first. The formulation is what is being offered.

Where several entrant sets are equilibria, the assortative one is the
cheapest, and under an anonymous gain schedule it therefore maximises the
unweighted sum of payoffs (Proposition~\ref{prop:assort}) --- so using it is
a stated selection with an argument behind it, not an assumption. That
argument is not a welfare claim, for the reasons set out with the
proposition.

Two features make the claim well-posed, and both required proof. The object
being signed is an invariant of the game: every pure-strategy equilibrium has
the same entrant count (Proposition~\ref{prop:count-inv}), so ``how many
compete'' does not depend on which equilibrium is selected --- even though
\emph{who} competes does (\S\ref{sec:primitives}, Information). And the class
of spreads for which the rule is proved is not bespoke: a pivot-spread is
exactly a dispersive-order shift with a crossing point
(\S\ref{thm:P9gen}), the dispersive order being the inequality order this
literature already uses \citep{HopkinsKornienko2009,Shaked1982}.

\paragraph{Why the extensive margin is the point.} The intensive-margin
version of the same question exists and is \emph{not} sign-determinate.
\citet{SchroyenTreich2016} sign a mean-preserving spread in wealth inside the
same privilege-contest payoff: their Theorem~3 gives the effect as positive
iff $2A(1-m^2) > P$, with $A = -u''/u'$ and $P = -u'''/u''$ --- local, two
players, and conditioned on the contest technology and the \emph{third}
derivative of utility. The dependence on $u'''$ is not cosmetic: at $m=1/2$,
CARA with $a=1$ and log utility at $w=1$ share $A=1$ but have $P=1$ and
$P=2$, and give opposite signs. The extensive-margin rule of
Theorem~\ref{thm:P9gen} needs only that $\kappa$ is strictly decreasing,
which is $u''<0$. So the margin is not a modelling convenience; it is what
makes inequality's effect on competition determinate at all. An anticipated
objection should be met here rather than left to a referee: the two results
answer the same economic question --- what does a spread of the resource
distribution do to competitive expenditure --- in different models, and the
comparison of hypotheses across margins is legitimate precisely because the
margin is the object of choice being defended. It is also not robust
terrain for the intensive side: in the assignment literature the direction
of the inequality effect reverses when the allocation re-optimises
(\citet{CostrellLoury2004}, Proposition~10 against their Proposition~6), so
conditioning results of this kind on the margin is already the literature's
own practice.

\paragraph{The 1981 precedent.} \citet{LazearRosen1981}, in the ``Income
Distributions'' subsection of their \S III, already have agents self-selecting
into a tournament by endowed wealth. This must be acknowledged, with two
qualifications that bound what it concedes.

First, their claim is offered as an example and they say so: the subsection
opens ``While it is not possible to make a general argument based on an
example, table~1 suggests\ldots'', and rests on one utility function, quadratic
costs, normal errors and two endowment levels. They \emph{suggest} the sorting
under declining absolute risk aversion; the present model \emph{derives} it,
distribution-free.

Second, the channel differs, and decisively. Their sorting runs through DARA,
which is a restriction on $u'''$: $A'(w)<0$ has numerator $(u'')^2 - u'''u'$
and so requires $u'''>0$. Here the sorting runs through $\kappa(w) =
u(w)-u(w-c)$ being decreasing, which follows from $u''<0$ alone. Quadratic
utility separates them: it is concave, has $u'''=0$, and exhibits
\emph{increasing} absolute risk aversion --- so $\kappa$ still sorts while DARA
fails outright.

Third, and least often noticed, the object of choice differs. Theirs is a
choice between two compensation schemes that both pay, so the sorting is
tolerance for the binomial spread of a lottery against a concentrated
alternative. Ours is whether to pay a fixed entry cost at all, so the sorting
is affordability in utility units. That is a difference in the economics, not
only in a derivative.

\paragraph{What is not claimed.} The mechanism --- wealth sorting entry
through concavity --- is not novel; the two precedents above own it between
them, and the sentence claimed here is the two-branch sign rule, not the
sorting. The remaining results are apparatus in its service and are
presented as such. The uniform cap of \S\ref{prop:P7} is the deterministic,
heterogeneous, utility-unit form of an accounting bound already present in
\citet{FuJiaoLu2015} (equation~(2) with Definition~2) and \citet{FuLu2010}
(equation~(7)); it is used here to dispose of the ceiling objection, not
claimed as new. The characterisation of \S\ref{prop:PMU} runs on the
order-statistic kernel of \citet{RyvkinDrugov2020}, whose message --- that no
competitor-number prediction is universal across noise distributions --- is
theirs in print; what is added is the extensive-margin instance that refutes
the referee's conjecture. And the scope of the central claim is stated with
it. The rule is proved under a one-sided condition on the quantile difference
down to the margin (Proposition~\ref{prop:tailprop}); pivot-spreads ---
equivalently, dispersive shifts with a crossing --- are the special case in
which that condition holds at every margin. A mean-preserving spread is
covered when, and only when, its quantile difference is signed down to the
margin; it may cross freely below. That is a restriction and not a
formality --- mean-preservation does not imply it, and spreads violating it
can move $k^*$ the other way (\S\ref{prop:tail}). What is not settled is a
characterisation of which spreads have the property (\S\ref{sec:open},
item~\ref{item:tailopen}).

\section{How this addresses the referee reports}

\subsection{Asymmetric \texorpdfstring{$k$}{k}-of-\texorpdfstring{$Q$}{Q} equilibria}
\label{sec:R1-asymmetric}
Both referees identified the same defect: the prior version compared
aggregate participation regimes rather than checking unilateral
deviations, which is not Nash logic for a finite-player game.
Proposition~\ref{prop:P5} (\S\ref{prop:P5}) replaces the binary
all-in/all-out structure with the asymmetric $k$-of-$Q$ equilibrium
requested, derived from first principles rather than assumed.

\subsection{Prize versus affordability}
\label{sec:R2-affordability}
Referee 2 objected that, with $y_2-y_1$ serving simultaneously as
``inequality'' and the contest prize, the claim that inequality raises
status expenditure was close to tautological, and that a rise in the
prize should raise \emph{all} participants' incentives, not only the
incumbent's. Decoupling $V$ from the wealth distribution
(\S\ref{sec:primitives}) and separating the two channels in
Proposition~\ref{prop:P8} (\S\ref{prop:P8}) addresses this directly;
Theorem~\ref{thm:P9} (\S\ref{prop:P9}) then states the
inequality result as a genuine theorem rather than a definitional
artefact.

\subsection{The mobility ceiling}
\label{sec:ceiling}
Both referees identified the old ceiling as an artefact: mechanical
arithmetic under uniform supports, a $Q\to\infty$ limit inconsistent
with the equilibrium existing, and a fixed prize with unbounded entry
not corresponding to anything. Proposition~\ref{prop:P7}
(\S\ref{prop:P7}) replaces it with the saturation bound $k^*\le
V/\kappa$, uniform in $Q$ and requiring no limit to be taken.

\section{Open items}
\label{sec:open}

Closed since the previous draft: the P9 fixed-point effect (there is no
feedback to resolve --- Proposition~\ref{prop:anon}); uniqueness of the
wealth fixed point (Lemma~\ref{lem:crossing}); the characterisation of
$\mu^*$ (\S\ref{prop:PMU}); P8, now proved rather than verified;
non-linear spreads (Theorem~\ref{thm:P9gen}); and P9 strictness
(Proposition~\ref{prop:strict}). Also closed: the literature position, which
was previously listed here as gating further proof effort. It is settled in the
accompanying literature survey and stated in
Section~\ref{sec:contribution}; the verdict is that the mechanism has
precedents and the contribution is the combination of P7, P9-gen and P-MU.

What genuinely remains:

\begin{enumerate}
\item \textbf{The incumbent's own entry decision} is exogenous
  throughout ($C$ fixed in $H_m=F^mG^{Q-1-m}C$). Endogenising it, and
  testing whether an equilibrium exists in which the incumbent
  abstains while a challenger invests, is untouched. This is a
  modelling extension rather than a gap in what is claimed.
\item \textbf{Spreads that are not pivot-spreads --- resolved, see
  Proposition~\ref{prop:tail}.} This item previously read that the sign rule
  was open for general Rothschild--Stiglitz spreads, and proposed importing
  the quantile-decomposition machinery of \citet{CostrellLoury2004} and
  \citet{Suen2007}. That turned out to be unnecessary. Those techniques sign
  a smooth \emph{aggregate} under a general spread and therefore need
  integration by parts and, in Suen's case, log-concavity of $1-F$. But
  $k^*$ is not an aggregate: it is a threshold statistic that reads the
  quantile function only down to the margin. The recorded ``count versus
  aggregate'' obstacle was real, and it runs the opposite way --- the count
  is \emph{easier}. What remains genuinely open is narrower.

\item \textbf{When does the tail condition hold? --- resolved, and the
  answer is that the hypothesis is sharp.}
  \label{item:tailopen}
  This item previously asked for a characterisation of the mean-preserving
  spreads satisfying Proposition~\ref{prop:tailprop}. That characterisation
  turns out to be exact but nearly definitional, and the substantive question
  --- whether the hypothesis could be weakened --- has a negative answer. Both
  are established in \S\ref{prop:band} and
  \file{checks/verify\_tailband.py}. What remains open is narrower and is
  stated there: whether some statistic of the spread \emph{other} than the
  displacement signs predicts the direction inside the band. That is where the
  \citet{CostrellLoury2004} monotone-weight technique would belong if an
  aggregate version were wanted later.
\item \textbf{Strictness away from the linear family.}
  Proposition~\ref{prop:strict} is stated for $T_\lambda$ with
  $\lambda\to\infty$. A characterisation of exactly which spreads move
  $k^*$ strictly, rather than an asymptotic sufficient condition, is not
  attempted.
\item \textbf{The competitor-number effect, sharpened.} $\Delta(0,\cdot)$ is
  quasi-concave in $Q$ when $\varphi$ is single-peaked, and $\varphi$ is
  single-peaked when $r$ is uniform (\S\ref{prop:PMU}). Open are the laws of
  $r$ beyond the uniform for which $\varphi$ is single-peaked, and the location
  of the peak in $Q$. An earlier version of this item conjectured an interior
  \emph{minimum}, from the reversed orientation that \S\ref{prop:PMU} corrects.
  Under the single-peaked hypothesis an interior minimum is impossible.
\end{enumerate}

\bibliography{refs}

\end{document}
